\begin{center}
{\Large\bfseries Selected Hecke modules and the denominator estimate}\\[4pt]
Edition 0.84 \quad 25 November 2026
\end{center}

The global construction selects a family of arithmetic data. To apply
the denominator estimate one must use the same member, its integral
Hecke modules and its chosen local characters. We prove that two
localisations above a common place and row index coincide whenever
the same good row satisfies the prescribed congruences at both.
The proof identifies their residual maximal ideals inside the Hecke
algebra acting on the selected module.

The finite central operator picks out the good part of the theta
spectrum. The resulting cut preserves the required Galois and
coefficient data under the original Galois preconditions.
The alternative denominator hypotheses give the estimates on one
hub with the scale, norm choice and coefficient comparisons required
by the global construction. At a type prime, one finite normal
extension makes every retained representation crystalline.
On the period spaces the actual inertia characters and the trivial
character form a Frobenius-stable set of size at most $n_V+1$.
A calculation of the integral lifting exponent identifies the three
monodromy estimates still needed there. We use the local data in a
bound for finite cocycles and relate
the integral congruence ideal to the remaining global estimate
on the selected family.

\begingroup\emergencystretch=3em
\section{The original compatibility conditions at a selected member}

% BEGIN M84T1
\paragraph{M84T1. The period reading already present at an original member.}
Fix the local correspondence $h_L$ in the source and an actual
member in $\mathrm{MemberAtLayerPoint}(h_L)$. Write $D=w.D(Z)$.
For every retained coefficient index $i$ and $v\mid\ell_i$,
the inertia polynomial of $D.\mathrm{wdOf}_i(r,v)$ for a continuous
integral representation $r$, potentially crystalline at $v$, is the
inertia characteristic polynomial of its potentially semistable
Weil--Deligne realization, using the very period bridge $D.\mathrm{nbr}$
in the test. Here and below the coefficients are mapped by the selected
embedding $\iota_i:\overline{\mathbb Q}_{\ell_i}\simeq\mathbb C$.

To identify the input on this member, unpack its $c.a.\mathrm{nt}$
realization, its equality of data and its character-tuple equality.
The three projections used here are the Weil--Deligne reading
$A.\mathrm{wdT}$, the global-to-local comparison $A.\mathrm{wd}$,
and the equality $\mathrm{Au}.\mathrm{wdOf}$ at the added localization.
They are the three fields copied in the construction of a
$\mathrm{CharacterChoiceData}$ from this realization. They can be
read before constructing its other fields. The last equality identifies
the native parameter with the one read by $A.\mathrm{wd}$.
The inertia clause of $A.\mathrm{wdT}$ identifies its inertia
polynomial. The above-$\ell$ clause of $A.\mathrm{wd}$ reads it
on a bridge inducing the specified place above $v$.

Here is the comparison with the particular bridge in the test.
Any two embeddings of the chosen global algebraic closure into
$\overline{\mathbb Q}_\ell$ inducing that place differ, on their
images, by a local Galois conjugacy after identifying the two
base completions over the dense number field. Indeed the resulting
isometry of the dense algebraic subfields extends to the completions;
over the common completed base the two embeddings of its algebraic
closure are conjugate, by extension of embeddings. Transport the
Galois action on period tensors by this conjugacy, acting on the
period factor, with the conjugate restriction of $r$ on the vector
factor. This identifies the invariant modules: the subgroup fixing
a vector is carried to its conjugate. It commutes with period
Frobenius and monodromy and conjugates the two actions of each
original inertia element. In particular potential semistability
and the characteristic polynomials agree. To compare finite
coefficient models, extend both to their compositum. The matrices
of the common representation then agree entry by entry under the
fixed embeddings in $\overline{\mathbb Q}_\ell$. Extension of
coefficients commutes with taking period invariants here: a basis
for the finite coefficient extension gives a finite direct sum of
the original invariant equations. The induced scalarized period
spaces and their inertia actions thus give the same characteristic
polynomials. This argument preserves the complex embedding and
the specified inertia element while changing the auxiliary
period realization.
% END M84T1

% BEGIN M84T2
\paragraph{M84T2. The inertia polynomial kills every specified crystalline period.}
On every original member of M84T1 the full third condition in
$\mathrm{Tie4}$, namely $\mathrm{TieInePoly}$, holds. Fix the
index, place, normal finite-index local extension in the test,
continuous rank-$m$ representation, inertia element $s$, and
polynomial $q\in O_i[X]$ whose coefficient image is the native
inertia polynomial. Let $d$ be any of the tested potentially
crystalline periods. The assertion is
\[
  \sum_{j=0}^{\deg q} \mathrm{toLL}(q_j)\,
                      \mathrm{act}_{V,B}(e(s^j))d=0 .
\]
Here $e$ is the original bridge map on the decomposition group.
Potentially crystalline periods are potentially semistable under
$B_{\rm cris}\hookrightarrow B_{\rm st}$. After a finite coefficient
extension the original model and the one giving the Weil--Deligne
reading agree by M84T1. Write $A_s$ for the action of $s$ on its
$D_{\rm pst}$. Inertia fixes the maximal unramified scalar field,
so $A_s$ is linear there. It also commutes with period Frobenius
$\varphi$. In degree zero of the Weil homomorphism the Frobenius
correction and the unipotent factor used for Frobenius
semisimplification both act as the identity. Therefore the raw
inertia action on the complex period space has characteristic
polynomial $\iota_i(q)$, by M84T1. Cayley--Hamilton gives
\[
             \mathrm{wdMk}\big(q(A_s)y\big)=0
                    \qquad (y\in D_{\rm pst}).
\]

For clarity this last map evaluates a single coefficient component.
Over $\mathbb Q_\ell^{\rm ur}$ the scalar algebra of a finite
coefficient field $E'$ is $E'\otimes_{\mathbb Q_\ell}
\mathbb Q_\ell^{\rm ur}$. Its field factors are indexed by the
embeddings of the maximal unramified subfield of $E'$: the remaining
totally ramified extension stays a field over the maximal unramified
extension, by its Eisenstein uniformizer polynomial. Scalarization
is faithful on its chosen factor since it extends a field embedding.
Frobenius permutes the factors transitively and is bijective on
$D_{\rm pst}$. As in OFL7, applying the last display to every
$\varphi^j y$ detects every component of $q(A_s)y$, since
$\varphi$ commutes with $A_s$ and fixes $E'$ acting as coefficient
scalars. Hence $q(A_s)=0$ on the whole period module.

The coefficient inclusion on $V$ is injective, and remains so
after tensoring over $\mathbb Q_\ell$ with the period ring.
Likewise the inclusion of crystalline in semistable periods is
injective, since the vector space is flat. A period fixed over
the given finite extension maps into $D_{\rm pst}$, and both
inclusions intertwine $s$ and the coefficient scalars. Pulling the
polynomial identity back gives $q(\mathrm{act}(e(s)))d=0$.
On pure tensors $\mathrm{act}(e(s))^j=\mathrm{act}(e(s^j))$
by the group laws on the two factors. Expanding $q$ proves the
display with the native coefficients and every extension, rank,
matrix representation, polynomial and period in its quantifiers.
% END M84T2

% BEGIN M84T3
\paragraph{M84T3. The central action and the full Hecke eigenspace.}
The following constructs the integral-action and eigenspace parts of
the original member's sought $\mathrm{Pend}$ witness from its
localized type sections and integral class-sum test. At index $i$
put $M=D.\mathrm{Mt}_i$, $V=K\otimes_O M$, where $O=O_i$ and
$K=\operatorname{Frac}(O)$. Let $C$ be the complex space of forms,
$C_x$ the sum of the type spaces of representatives of $x$, and
$C_\emptyset$ the corresponding sum for representatives with no
Hecke point. Use the original independent type spaces and their
decomposition in the row. Suppose $\Phi:V\to C$ is a localized
type section: it is injective and semilinear for the given
embedding, $\Phi(V_x)\subset C_x$, and the difference between a
rational form and $\Phi$ of its projection is in $C_\emptyset$.
Assume the finite group action preserves each type space and
integral central elements send each distinguished rational lattice
generator into the lattice of rational forms. These are precisely
the section, type stability, and $\mathrm{CenInt}$ inputs. Use
the native equality identifying $M$ with the $O$-span of the
projections of these generators.

The separating elements of LN3 give $V=\bigoplus_x V_x$.
Thus $\Phi(V)\subset\sum_x C_x$. On a rational form $y$, its
projection has $\Phi$-image equal to the component of $y$ along
$C_\emptyset$ in $\sum_x C_x$. This projection commutes with
each group-algebra action because both summands are stable. For
an integral conjugacy class sum $z$, apply the integral test to
each distinguished generator. Projecting the rational lattice
combination it gives shows that $z\Phi(1\otimes M)$ belongs to
$\Phi(1\otimes M)$: algebraic integers of the global coefficient
field map into $O$ at this place by the defining valuation.
Every central element of $O[G_{0g}]$ is an $O$-linear combination
of these integral class sums, since conjugation permutes the
group basis. Hence the full center preserves $\Phi(1\otimes M)$.
Pull its action back by the injective $m\mapsto\Phi(1\otimes m)$.
This defines an $O$-algebra action integrally, even if $\ell_i$
divides the group order. Rationalizing recovers the given action
on the image of $\Phi$ by semilinearity. It commutes with Hecke:
the central operator preserves each $C_x$, and consequently every
$V_x$ since $\Phi$ is injective and the different $C_x$ are
independent. Hecke acts by the specified point character on that
space. For any rational forms $a,b$ whose complex values are
related by the central action, project their equality along
$C_\emptyset$ to get the required equality of their projections
in $V$. If a rational form is in a single type space this same
projection lies in the indicated eigenspace or vanishes for an
absent point. Finally, if a central element acts by the chosen
$a\in O$ on every type space above $x$, it acts by its image on
all of $C_x$. Applying the injectivity of $\Phi$ gives the scalar
$a$ on the entire $V_x$. This is the original $\mathrm{EigSpan}$
quantifier. The construction costs no denominator from a group
order or a point projector.
% END M84T3

% BEGIN M84T4
\paragraph{M84T4. Rational and integral central operators in the selected row.}
Both the rational test $\mathrm{ZLStable}$ and the integral test
$\mathrm{CenInt}$ in M84T3 hold on the actual members. The
hypothesis of membership includes an $\mathrm{RXs}(t)$
witness. Its $\mathrm{AllT}$ realization identifies the finite
group $G$ with $G_{0g}$, its level-coordinate spaces with the
rational forms at $j=k+1$, and the rational form lattices with
the specified generator spans. For a central group-algebra
element over the global coefficient field, pull its coefficients
back by the group isomorphism. The resulting central operator
on level coordinates is defined over that field, and
$\mathrm{AllT}.\mathrm{cen}$ identifies its complex image with
the operator on forms. Applying it to the coordinates of any
rational form proves $\mathrm{ZLStable}$.

For the integral assertion let $E$ denote the inclusion matrix
from level coordinates into the ambient form space, and let
$L$ be the ambient integral lattice denoted by $\mathrm{nabLat}$.
The level lattice is $\{x:Ex\in L\}$. Right translations by the
finite group preserve $L$ in these row models: in the function
description the translated value at a class is obtained from
the value at its translated class by the weight action, which
carries the value lattice there onto the value lattice here.
This is the $\mathrm{nab}$ component of the integral action
test; the same change of class preserves the row Hermitian
pairing, which is positive definite at the chosen infinite
place. For $a=\sum_g n_g g$ with integral coefficients and $a$
central, $\sum_g n_g r(g)$ commutes with the idempotent imposing
the level character. Its image therefore preserves the level
subspace, and its level matrix $a_{\rm lev}$ satisfies
\[
       E a_{\rm lev}x=\sum_g n_g r(g)Ex\in L\qquad(Ex\in L).
\]
This proves integral preservation on level coordinates. Under
$\mathrm{AllT}.\mathrm{latL}$ each distinguished generator
belongs to this lattice (take its own coefficient to be one),
and every lattice output is a combination of the distinguished
complex generators over the global integers. Lift those
generators back by the rational coordinate isomorphism and
take the same combination. Its complex image equals the
central translate of the original generator by the operator
identity above. This is exactly $\mathrm{CenInt}$ in every
retained row.
% END M84T4

% BEGIN M84T5
\paragraph{M84T5. The two coefficient fields for the series idempotents.}
In the eventual quantifier defining membership take the fourth
finite list of prescribed algebraic complex numbers to contain
$\exp(2\pi i/N)$, where $N>0$ is the base's finite type-group
order. Its algebraicity follows from $X^N-1$. This list depends
only on the permitted base data, before $w,Z$. The global
coefficient field $L$ then contains a primitive $N$th root by
the defining inclusion test and injectivity of its complex
embedding. This is an allowed strengthening of the eventual
threshold. In the fraction field of every $O_i$ its image is
again primitive and is integral over $O_i$ by $X^N-1$. Since
$O_i$ is integrally closed, it lies in $O_i$ with the same
order. This supplies the exact $\mathrm{zetaPrimitiveRoot}$
entry for the desired original witness.

For either type prime $g_k$, $|\mathrm{GL}_r(\mathbb F_{g_k})|$
divides $N$ by the product formula for the type-group order:
the corresponding group order occurs to the positive power
$D.R.f$. In a complex irreducible representation $\tau$ of
that finite group every eigenvalue of $\tau(h^{-1})$ is an
$N$th root of unity. Thus its trace belongs to the complex
image of $L$. The coefficient at $h$ of the irreducible central
idempotent is
\[
      \frac{\dim\tau}{|\mathrm{GL}_r(\mathbb F_{g_k})|}
                          \operatorname{tr}\tau(h^{-1}).
\]
The member's $A.\mathrm{grp}$.series identity expresses the
specified idempotent as a sum of such idempotents. Its
coefficients therefore belong to the image of $L$. This proves
$\mathrm{SerLL}$ with the native coefficient field. Dividing
by the group order in this computation takes place in $L$.
% END M84T5

% BEGIN M84T6
\paragraph{M84T6. Five original Galois preconditions on the same member.}
The fields $\mathrm{large}$, $\mathrm{nct}$, $\mathrm{same}$,
$\mathrm{block}$ and $\mathrm{muUnr}$ of $\mathrm{GaloisPreconditions}(D,n)$ hold.
The coefficient condition is the $\mathrm{large}$ part of the
member's given nonempty $\mathrm{defaultExtension}$. The
norm-choice tie is the tie of the selected norm datum.

For $\mathrm{same}$, fix a localization $a$ over $w$. Its
actual $\mathrm{NdSpec}.\mathrm{tw\_reads}$ gives the
native homomorphism $f$ out of the initial coefficient ring
and embeddings of the two fraction fields into a field
$\Omega$. In $\Omega$ their readings agree on every element of
the initial ring and on every element of the global coefficient
field. The norm-choice tie gives a single global character
$\psi_0(w)$ whose evaluations at $g$ read as the two specified
local character values (with the prescribed Artin map on the
norm-choice side). Thus the images in $\Omega$ of every finite
sum $\sum_g f(a_g)\psi(g)$ are equal for the two readings.
The ring maps from the local rings into $\Omega$ are injective:
the small-to-large norm-choice ring map is injective, and all
remaining maps factor through inclusions into fraction fields
and field embeddings. Hence each such finite sum vanishes
at one reading exactly when it vanishes at the other. This
is $\mathrm{SameChar}$, and the index comparison in the
norm-choice tie gives $\mathrm{SameTwist}$ for each $a$.

For $\mathrm{block}$ put $j=D.\mathrm{nd}.j(i)-1$ as a
zero-based index and $e_\rho=(D.\mathrm{hc}\,\rho).e(j)$.
The member already carries the automorphic character tie
$c.a.\mathrm{chrN}$. It identifies
$\mathrm{tw}(\mu_C(i),\eta_T^{-1})$ with the tuple's
$j$th character times $\eta^{-1}$. The finite-order multiplier
of the tuple is trivial on the connected archimedean
components. Its infinity exponent is therefore that of the
underlying $\mu_0$, namely $-e_\rho$ in the chosen CM type.
The exponent of $\eta$ on that type is $1/2$. Inversion
negates exponents and multiplication adds them, as follows
directly by multiplying the character formulas on the
archimedean units. The two conjugate embeddings have opposite
angle exponents. The infinity-parameter clause of the same
character tie gives precisely the two singleton multisets
$\{-e_\rho-1/2\}$ and $\{e_\rho+1/2\}$ in the designated
order. These are the original $\mathrm{BlockCharParam}$ tests.

At a place $w$ above $\ell_i$ outside the two generator
primes, the same twisted rank-one automorphic parameter is
unramified. Use the unramified-place equivalence in
$c.a.\mathrm{chrN}$ to check the corresponding Hecke
character. The character $\eta$ is unramified here, and
$\mu_0$ can ramify only above a generator prime. In the tuple
multiplier $\delta\psi$, the chosen $\delta$ can ramify only
at its designated auxiliary prime. The norm choice uses the
other auxiliary character when $\ell_i$ equals the first
prime; the auxiliary primes are distinct. The character
$\psi$ is unramified outside its selected $Q$ at such a
place. The localization identifies this $Q$ with $Q_i$ by
$\mathrm{NdSpec}.Q_{\rm eq}$ and its base-$Q$ identity, and
the defining $S_i,Q_i$ condition excludes places above
$\ell_i$ from $Q_i$. Each factor is therefore trivial on
local units, as is their product. This proves
$\mathrm{MuUnr}$ with its coefficient-prime guard.
% END M84T6

% BEGIN M84T7
\paragraph{M84T7. Localized type sections from the member's Hecke realization.}
The section $\Phi$ used in M84T3 exists on the original member at
each $i$. Here are the Hecke inputs actually needed. At a split
place outside the excluded set, the member's $\mathrm{MtInput}$
gives hyperspecial level and trivial character there. Indeed
the two type-prime places are excluded, while the ramification
of the tuple multiplier is above the norm choice's $Q$ or the
residue characteristic, both excluded by the localization.
The chosen $\mathrm{NdSpec}$ has a rational operator on forms
realizing each of the specified spherical double cosets. This
also follows by taking $\mathrm{MtInput}.\mathrm{rat}$ in the
row $j-1$. Its projection is the specified local Hecke
operator. To see this, use the injective comparison $\Psi$
of the rationalized local lattice with the part of the
normalized Hecke space supported at its maximal ideal. On
the whole space the rational operator is the corresponding
generator, by its definition on rational forms. Its action
on the retained part is the localized action by
$\mathrm{NdSpec}.\mathrm{act\_hK}$. The projected rational
form differs from the original by a vector in the other
parts by $\Psi_{\rm pr}$. Those parts are stable under the
commuting Hecke generators. Thus the two candidate answers
differ by a vector in both the retained and the complementary
parts. These are disjoint by decomposition into the spaces
of distinct simultaneous characters. Injectivity of $\Psi$
gives the assertion.

The member's automorphic specialization gives precisely its
irreducible type spaces and their point test: equality of all
split-good Hecke scalars with those of $x$ is equivalent to
having point $x$. On each such type space each tested operator
is scalar. Indeed, at a split place the fixed space of a
hyperspecial subgroup in an irreducible spherical factor of
$\mathrm{GL}_{n_V}$ has dimension one, and the Hecke sum acts
only at this factor. All copies of the same representation
therefore have the same scalar. Equivalently, without choosing
a tensor factorization, the fixed space in the global
irreducible representation is a simple module for the
spherical Hecke algebra and the translations away from the
place. The translates of any nonzero vector generate the
irreducible representation, and averaging a finite sum of
such translates over the hyperspecial subgroup recovers the
fixed vectors using spherical double cosets. Spherical
convolution commutes (transpose preserves every double coset
by elementary divisors and reverses convolution), and it
commutes with the away translations. Its action on this
simple module is scalar: smoothness and second countability
bound the dimension of a cyclic representation by countable
dimension. A nonscalar endomorphism $T$ in its division
endomorphism ring would make $(T-a)^{-1}y$, $a\in\mathbb C$,
linearly independent for any $y\ne0$: multiply a finite
relation by the common denominator and evaluate the resulting
polynomial at each $a$. A nonzero polynomial in $T$ is
invertible since it splits into linear factors. This contradicts
countable dimension. Finally the level and hyperspecial Hecke
sums agree on these fixed vectors, since away from the place
one has the same support and at the place the same coset
representatives. The normalized scalar is the same on every
copy under a Hecke intertwiner.

Choose a basis of each $V_x$ in the direct sum of M84T3, with
coordinates $\lambda_{x,b}$. Each functional
$y\mapsto\iota_i\lambda_{x,b}(\mathrm{pr}\,y)$ on rational
forms extends to a complex-linear functional on $C$ by the
member's rational basis of forms, whose complex images are a
basis. Write $p:C\to\mathbb C^{\sum_x\dim V_x}$ for their
joint map. If the $x,b$ coordinate is nonzero on a vector
in a representative's type space, apply each of the tested
operators to that vector. On rational basis elements the
extended functional intertwines that operator with the
eigenvalue at $x$, by the projection just proved, hence it
intertwines on every vector. Cancelling the nonzero coordinate
identifies the representative's scalar with the scalar at $x$
for every split test. The point test then gives $x$ as its
point. Hence $p$ takes $C_x$ into its own coordinate block
and kills $C_\emptyset$.

The projections of rational forms span $V$ over $K$.
Consequently the image of $p$ contains the coordinate
basis, and $p$ is surjective. Since $C$ is spanned by the
type spaces, its restriction to $G=\sum_x C_x$ is still
surjective. The member's multiplicity identity gives
$\dim C_x=\dim_K V_x$, so $\dim G\le\sum_x\dim_K V_x$.
Thus $p|_G$ is an isomorphism. Take its inverse on the
embedded $K$-coordinates of $V$. The resulting map is
injective, semilinear, and takes the $x$-space into $C_x$
(the block images already span the individual blocks).
Its composition with projection on rational forms differs
from the original complex image by $\ker p=C_\emptyset$.
These are exactly all fields of $\mathrm{Sec}$.

In particular the remaining hypothesis in M84T3 is the
preservation of the individual type spaces by the action.
It suffices to identify the complex central operator on the
actual realization with its sum of adelic right translates
(the member comparison $\mathrm{CenLoc}$): each isotypic
sum in this realization is right-translation stable. On
noncentral elements the native $\mathrm{onMt}$ is zero.
% END M84T7

% BEGIN M84T8
\paragraph{M84T8. From local parahoric vectors to the fourth original tie.}
Fix a series system $S_k$ for each type prime whose rank-$r$
projectors are the original $e_{\rm Ser}(f)$. Assume the
local representation comparison $\mathrm{LocType}(D,S)$
at both primes. It states that the nonzero central actions
on type vectors have parahoric block vectors in a single
local representation $\pi$, supplies the determinant-character
label on the lower block and the matching inertia value of
$\mu_C$, and identifies the inertia polynomial of the
automorphic parameter from any full-rank series detected in
$\pi$. The member satisfies the additional unramified-twist
identity $\mathrm{IneUnr}(D)$: its $A.\mathrm{wdT}$ carries
rank-one twist to tensor product and preserves unramifiedness,
dimension and inertia polynomials. On every inertia element
the unramified rank-one factor acts identically, so tensoring
it with any matrix preserves the characteristic polynomial.
Thus the full $\mathrm{SerLoc}(D)$ follows from the stated
series system and local representation comparison.

Fix a parameter $\sigma$ for which $e_{\rm Ser}(f)$ acts
nontrivially on a vector in its parameter space. The native
span identity for the parameter space then gives a single
type space on which this action is nonzero: otherwise its
linear kernel would contain their span. Choose its actual
representative $P$ and the list $\Sigma$ in the span test.
Its isobaric localizations give $\sigma$, and appending the
distinguished rank-one string and twisting by $\eta_T^{-1}$
is a base change of $P$. The local comparison gives $\pi$,
$\chi,c$, and a parahoric vector with nonzero top projection.
In the principal-congruence fixed space write $u$ for that
projected vector. It is fixed by the parabolic unipotent
radical: translating by a top-block matrix conjugates the
radical into itself. The top projector fixes $u$ by
idempotence. The lower-block projector for
$h=(X-c)^t$ acts on the original vector by a nonzero
idempotent scalar on the determinant character $\chi$,
hence by one; it commutes with the top action, so fixes
$u$ as well. The product property in the series system
therefore gives $e_{n_V}(fh)u=u\ne0$.

Write $[f]_\tau$ for the product of $X-\tau(x)$ over
the roots of $f$ in an algebraic closure of the residue
field, with multiplicities, where $\tau$ is the specified
injective Teichm\"uller map. The full-rank comparison gives
$[fh]_\tau$ for the inertia polynomial on $P$ at the
specified inertia generator. On the base-change list this
same polynomial is
\[
       [\mathrm{labPar}(\sigma)]_\tau\,(X-\tau(c))^t.
\]
Indeed sums multiply inertia polynomials, the empty sum has
polynomial one, and the real unramified shifts in a Speh
string have trivial inertia, so a string of length $d$
contributes the $d$th power. The rank-one $\eta_T^{-1}$
twist disappears by the unramified-twist test. The
degree of the parameter's label equals its rank, so the
native degree guard in reading its inertia polynomial
holds. The lower-block value is precisely the scalar
prescribed by $\mathrm{LocType}$ for $\mu_C$. Cancel
$(X-\tau(c))^t$ in $\mathbb C[X]$. The remaining
root multisets agree by injectivity of $\tau$, including
all possible repeated roots. The two polynomials are monic
(the first by the nonzero series projector and the second by
the parameter condition); factorization over the residue
algebraic closure and injectivity of base change now give
$f=\mathrm{labPar}(\sigma)$. This proves the asserted
series label for every vector in the original test. The
comparison $\mathrm{LocType}$ itself concerns irreducible
local representations and depth-zero types; M84T8 isolates
it as an input.
% END M84T8

% BEGIN M84T9
\paragraph{M84T9. The same tame inertia characters at all original type places.}
On the original member one can choose the entire
$\mathrm{om}$ family of the desired $\mathrm{Pend}$ witness,
including its eigenvalue and character-twist equalities.
The initial $A.\mathrm{lab}$ reads the Weil--Deligne label
at each selected generator $\gamma_{k,v}$. The given
$\mathrm{defaultExtension}$ supplies its conjugate-place
counterpart $\mathrm{labc}$, including a lift of CM conjugation
identifying the valuation rings and conjugating the generators.
At a selected place take the given generator, and at its
conjugate take the generator in this latter identification.
These assignments agree on their domains: above a real place
of the completely split type prime there are exactly two
places, exactly one in the chosen half. The selected-place
index is injective. Complete the assignment by the identity
away from the type prime. This proves the member's
$\mathrm{LabAll}$ and transports tame generation to every
conjugate place: the indicated conjugacy carries inertia,
its wild subgroup and all their open subgroups to the
corresponding groups on the other side.

For such a place $w$ over $g=g_k$, the residue field has
order $g$. The tame inertia quotient is procyclic of
prime-to-$g$ order, generated topologically by the chosen
$\gamma_w$. More explicitly over the maximal unramified
extension, adjoining $d$th roots of a uniformizer gives
the cyclic quotients of degree $d$ prime to $g$: all
$d$th roots of unity are already present by residue-field
lifting, and adjoining the root has degree $d$ by valuation.
These quotients exhaust the tame inertia action. For
$x\in\overline{\mathbb F}_g^\times$ take its finite order
$d$, which is prime to $g$. There is a unique tame character
of inertia with value $\tau(x)$ at $\gamma_w$: use the cyclic
quotient of order $d$, and uniqueness follows from tame
generation. Denote it by $\omega_{w,x}$, with $\tau$ the
specified multiplicative lift. Set the value to one off
inertia; use the trivial character for $x=0$ and at places
away from $g$. These assignments satisfy all the native
one, multiplicativity, wild, and openness tests. In detail
multiplication in $x$ follows by uniqueness on the generator,
and finite image gives an open kernel in inertia. The
induced profinite topology provides an open subgroup of the
global Galois group whose intersection with inertia lies in
that kernel. Distinct nonzero $x$ are distinguished already at
$\gamma_w$ by injectivity of $\tau$.

For any local parameter, the label records the eigenvalues
of this generator on its wild-invariant subspace, with
multiplicity. Indeed inertia on a complex Weil--Deligne
representation has finite image, and on the wild-invariant
part the generator has prime-to-$g$ order. Its eigenvalues
therefore have unique preimages under $\tau$; the multiset
is stable under $x\mapsto x^g$ by Frobenius conjugation
on tame inertia. It gives the unique monic polynomial over
$\mathbb F_g$ with that multiset of roots, exactly the
Weil--Deligne label in $\mathrm{LabAll}$. Its degree is
the dimension of this wild-invariant subspace. Under the
native full-degree guard the whole space is wild-invariant.
The finite tame action is a direct sum of eigenspaces of
the cyclic quotient. Its characters are uniquely determined
by the generator, hence are the $\omega_{w,x}$ just chosen
with the multiplicities of the label's roots. This proves
the entire polynomial identity for every inertia element.
It also covers the zero-dimensional representation, for
which the empty product is one.

Finally, suppose a rank-one parameter has label $X-u$ with
$u\in\mathbb F_g^\times$. It too has full degree, and its
inertia character is $\omega_{w,u}$. The member's
$A.\mathrm{wdT}$ identifies the rank-one twist with tensor
product on Weil--Deligne representations. Tensoring the
previous tame representation by this line stays tame of the
same dimension. Each root $x$ of the original label is
multiplied by $u$, because $\tau(ux)=\tau(u)\tau(x)$ and
the tame characters agree on their generator. This proves
the degree and multiset assertions for the twisted parameter.
Thus the very same choice supplies both
$\mathrm{inertiaCharpoly\_eigenvalues}$ and
$\mathrm{inertiaCharpoly\_character}$ on all original places.
The remaining local comparison in M84T8 and the rank-one
de Rham and depth-zero representation tests in the first
two clauses of $\mathrm{Tie4}$ have separate content.
% END M84T9

% BEGIN M84T10
\paragraph{M84T10. The first character tie on the original member.}
Fix an original coefficient index $i$, a generator prime $p=g_k=\ell_i$
and $v\in\mathrm{PlG}(k)$. Write $w_+=\mathrm{plGw}(k,v)$,
$w_-=\overline{w_+}$ and $j=j(i)-1$ in zero-based indexing.
Both completions are $\mathbb Q_p$ by complete splitting. We prove
$\mathrm{tieChi}$ using global and local Artin reciprocity, Chebotarev
density and the rank-one crystalline period calculation.

Let $\mu=\mathrm{tp.tup}(\mu_i,j)$, $\lambda=\mu\eta^{-1}$
and $s=(t-1)/2+(1-n_V)/2=-r/2$. At every split good prime $z$
outside the finite ramification set and away from $p$, the member's
$\mathrm{chrN}$ identifies the one Hecke eigenvalue of this line
with $\lambda(\varpi_z)^{-1}$. The Frobenius condition on
$\chi_1=\mathrm{nd.chi}(i,1)$ therefore reads
\[
 \iota_i\bigl(\chi_1(\mathrm{Frob}_z)\bigr)
        =\lambda(\varpi_z)^{-1}(\mathrm Nz)^s .
\]
Here the embedding on the left includes the indicated embedding of the
coefficient ring. Put $\beta=\lambda|\cdot|^s$, using the idelic
absolute value, so the right side is $\beta_z(\varpi_z)^{-1}$.
The angle exponents of M84T6 give the two infinity exponents of
$\beta$ above each real embedding as
$s-(e_\rho+1/2)$ and $s+(e_\rho+1/2)$. Both are integers because
$e_\rho\in\mathbb Z$ and $r$ is odd. Thus $\beta$ is an algebraic
Hecke character.

For precision take the inverse finite character $\beta_f^{-1}$ and
the Artin normalization sending uniformizers to the Frobenius in this
display. Its $p$-adic avatar can be read directly on finite ideles.
On a principal idele $x$, $\beta_f^{-1}(x)=\beta_\infty(x)$ is
the product of the indicated integral embedding powers. Hence all
values of $\beta_f^{-1}$ are algebraic and lie in a fixed number
field: the finite local unit characters have finite image, and
taking powers of finitely many ideal class representatives reduces
their values to principal values and unit values. Transport
$\beta_f^{-1}$ through $\iota_i^{-1}$ and divide by the product
of these same algebraic powers on the $p$-components, grouping the
embeddings by the places they induce. The result is trivial on
$E^\times$. It is continuous: at $p$ the correction is locally
algebraic and the finite part is locally constant, while away
from $p$ only finitely many unit groups act nontrivially.
It therefore factors through the profinite idele class quotient
by Artin reciprocity. Indeed modulo the closure of $E^\times$ in
the finite ideles the unit groups have compact image of finite
index, by finiteness of the ideal class group; the resulting
quotient is profinite. At a good $z\nmid p$ it takes the displayed
Frobenius value. These Frobenius elements determine a continuous
abelian character: given a finite Galois quotient over $E$, take
a Galois closure over $E^+$ and lift the desired element to the
subgroup fixing $E$. Chebotarev over $E^+$ gives a prime outside
the finite exclusions, splitting in $E$, and a place above it with
that Frobenius in the quotient. This proves equality of the avatar
with the original $\chi_1$.

At either $w=w_\pm$, the correction on local units is an integral
power of the identity under $E_w=\mathbb Q_p$. Local reciprocity
therefore expresses $\chi_1$ as a Tate power times a finite
character and an unramified character. Explicitly extend the smooth
unit character to the local multiplicative group by assigning one
to $p$, and remove this finite character and the indicated
cyclotomic power from $\chi_1$; the quotient is unramified.
The unramified line is crystalline by the rank-one period
calculation of S80J4; the Tate line has a period given by a power
of Fontaine's $t$. Their Weil--Deligne inertia is trivial.
After a finite extension killing the smooth character the remaining
finite line has the period $1\otimes e$. Inertia acts on this
potentially crystalline period through the original finite
character, since inertia fixes unramified period constants.
Finite-extension descent gives de Rhamness on $E_w$: after normal
closure the full de Rham periods descend by semilinear Galois
descent. The member's $\mathrm{loc\_tie}$ and M84T1 read this on
the required $\mathrm{nbr}$ and identify the native inertia
polynomial. It remains to identify the smooth character on one
specified generator at each place.

On units of $w_+$ the factors $|\cdot|^s$ and $\eta$ are trivial.
In $\mu=\mu^0_j\xi$ the multiplier is unramified there: its
selected $\delta$ can ramify only at its selected auxiliary prime,
distinct from the coefficient prime by the choice in M84T6, and
$\psi$ is trivial above that coefficient prime. This includes the
reference tuple with trivial multiplier. Thus the smooth unit
character is read from the tame character $\mu^0_j$, with the
fixed inverse from the avatar. Let $\gamma$ be the designated
tame generator and choose $\varpi\in w_+\setminus w_+^2$ in
$\mathcal O_E$. For $\pi^{p-1}=\varpi$, the residue of
$\gamma(\pi)/\pi$ belongs to $\mathbb F_p^\times$. Lift it to
$b\in\mathcal O_E$, necessarily a unit at $w_+$. This is exactly
the exponent-one test $\mathrm{TameRep}(w_+,1,\gamma,b)$.
The local Artin unit of $\gamma$ has residue $\bar b$ or
$\bar b^{-1}$ according to reciprocity convention. Indeed
$\pi/(\zeta_p-1)$ is an algebraic unit, whose residue inertia
fixes, and
$\gamma(\zeta_p-1)/(\zeta_p-1)$ reduces to the mod-$p$
cyclotomic character of $\gamma$. Local cyclotomic reciprocity on
$\mathbb Q_p^\times$ sends a unit to that unit or its inverse
on $p$-power roots of unity. The member's $A.\mathrm{chiT}$
now gives $\mu^0_{j,w_+}(b)=\tau(c)$, with
$c\in\{\chi_J,\chi_J^{-1}\}$. Since the character is tame, only
the residue of the Artin unit matters. Its value on $\gamma$ in
the Weil--Deligne line is consequently $\tau(u)$ for one
$u\in\{\chi_J,\chi_J^{-1}\}$, absorbing the fixed inversions.

Finally $\mu^c=\mu^{-1}$ on conjugate local units: apply the
$\mathrm{csd}$ equivalence in $\mathrm{chrN}$ to
$\mathrm{muC\_csd}$ and use its tuple equality, on finite ideles supported at the tested place. For the local
component first take $x\in E^\times$ as in the definition of
$\mathrm{csd}$, then extend to $E_w^\times$ by density.
M84T9 chooses the conjugate generator by transport, and local
reciprocity commutes with this transport. Thus the smooth
Weil--Deligne character at $w_-$ takes the value $\tau(u^{-1})$
there. Both inertia actions are tame on the full line; their native
labels are therefore $X-u$ and $X-u^{-1}$ by M84T9. These are
distinct places; for every other $w\mid p$ the two implications
in the definition are vacuous. This proves the first original
tie, with the same $u$ for both tests.
% END M84T10

\section{Localisation uniqueness on the selected integral Hecke module}

% BEGIN L84C1
\paragraph{L84C1. The character and row comparisons on an original member.}
Fix any selected member $D=w.D(Z)$, rows $t$, norm choice $n$ and
extension witness $c$ in the original member quantifiers.
The record $c.a.nt$ presents this very $D$ as the localisation of its
row and character data. Its Weil--Deligne identification at a place
$v$ preserves dimension, zero, sum, tensoring by a character,
real unramified twist, unramifiedness and Frobenius eigenvalues.
In the identified category a line has zero nilpotent operator.
Tensoring by a line distributes over a direct sum, commutes with
tensoring by the norm to a real power on either factor, and
successive character tensors compose by multiplication.
These isomorphisms intertwine both the Weil actions and the
nilpotent operators. The unramified line with Frobenius eigenvalue
one is the trivial line, so tensoring by it is the identity.
The bijection on isomorphism classes transports all five assertions
back to the member. These are the five laws of
\(\mathrm{RowTie.CharLaws}\).

At every split $v$ the two specialisations of a common row $\sigma$
at any $i,i'$ have local parameter equal to the one underlying
row parameter $\mathrm{locR}(\sigma,v)$. If $j_i=j_{i'}=j$ and
$w_i=w_{i'}=w$, the two tuple inputs $\mu_i,\mu_{i'}$ both equal
$\mu(w)$. In the selected character frame the two automorphic
characters $\mu_{C,i},\mu_{C,i'}$ both map to the $(j-1)$-indexed
entry of this tuple. Here $1\le j\le m$ makes this a common
zero-based entry. The frame identification is injective, so the
characters themselves agree. These comparisons prove the three
conjuncts of the built laws on the specified member.
% END L84C1

% BEGIN L84C2
\paragraph{L84C2. Comparing two membership polynomials before reduction.}
Suppose $i,i'$ have this same $(j,w)$ and $\sigma$ satisfies
$\mathrm{AtM}$ at both indices. Choose the two lists $\Sigma,\Sigma'$
and the two points $x,x'$ in their respective witnesses. At a common
split test place $v$, the lists separately have parameter
$\mathrm{locW}(i,\sigma,v)$ and $\mathrm{locW}(i',\sigma,v)$.
They are equal by L84C1. Appending $(1,\mu_{C,i},t)$ and
$(1,\mu_{C,i'},t)$ adds the same direct sum on both sides.
Indeed the term for a part $(a,\pi,d)$ is the sum of the local
parameters for $\pi$ with real twists $(d-1)/2-k$, $0\le k<d$.
Twisting the augmented list by $\eta_T^{-1}$ thus tensors the
entire sum by its local character. To see this including the empty
list, the zero local parameter tensored with this unramified
line still has dimension zero and empty Frobenius multiset;
the unramified parameter with that multiset is zero. Each added
summand uses the direct-sum and real-twist identities of L84C1
and compatibility of the local character with the global twist.
Finally apply the same real twist $(1-N)/2$, $N=n_V$, to both sides.
The resulting Galois polynomials over $\mathbb C$ agree.

Let $O_i,O_{i'}$ be the actual coefficient rings. The selected
localisation data identify both with $O_c$ for $(j,w)$, by
$\phi_i,\phi_{i'}$. The integral complex embeddings of the member
are the composites of these identifications with the common
embedding of $O_c$. Indeed the two given fraction-field embeddings
into the local algebraic closure and its complex embedding are
exactly the specialisations in the chosen data. The embedding on
$O_i$ is injective, as the composite of its map into its fraction
field and field embeddings. Consequently the polynomials $P_{i,v}$
and $P_{i',v}$ supplied by $\mathrm{AtM}$ have identical images in
$O_c[X]$. This is coefficientwise equality, before any use of the
Hecke congruence.
% END L84C2

% BEGIN L84C3
\paragraph{L84C3. The residual maximal ideal and the original localisation index.}
Fix the common $(j,w)$ of L84C2 and its selected Hecke algebra
$\mathcal H=O_c[H_{v,k}]$ inside the endomorphisms of the common
coefficient space. Its generators run over $\mathrm{GenAt}$.
At both indices they are split tests away from the residue
characteristic, the designated bad places and the common auxiliary
set $Q_c$. This uses $Q_i=Q_c=Q_{i'}$ and the equality of residue
characteristics forced by the common $w$. Under the selected
localisation map the scalar extension of the integral module
at $i$ injects into the common space with image the simultaneous
part at its maximal ideal $\mathfrak m_i$. Its point $x$
corresponds to an $O_c$-algebra eigencharacter
$\psi_x:\mathcal H\to O_c$ with nonzero eigenspace,
$\ker(\overline{\psi}_x)=\mathfrak m_i$, and
$\psi_x(H_{v,k})=\phi_i(h_{i,v,k}(x))$. The same description
applies to $i',x'$. All these identifications belong to the
selected normalised spectrum.

For every generator, compare the coefficient of $X^{N-k}$ in
the polynomials of L84C2. By the two congruences in $\mathrm{AtM}$,
after reduction in the common residue field it is on each side
the corresponding Hecke coefficient
\[
 (-1)^k q_v^{k(k-1)/2}\,\overline\psi_x(H_{v,k}),
 \qquad 1\le k\le N ,
\]
with $x'$ on the other side. This follows by expanding the monic
Hecke polynomial, since exactly the term with index $k$ has that
degree. Both coefficient isomorphisms respect the maximal ideals.
The scalar in front is a unit since $v\nmid\ell_i$. Thus the two
reduced eigencharacters agree on every generator and on $O_c$.
Every element of $\mathcal H$ is obtained by finite sums and
products of these elements, so the maps agree on $\mathcal H$.
Their kernels are equal. The map $i\mapsto\mathfrak m_i$ is injective
on the specified fibre of the selected normalised spectrum.
Therefore $i=i'$.

In the abstract Hecke-tied formulation the same argument uses the
given coefficient isomorphism compatible with complex embeddings.
If $i\ne i'$, its separation clause furnishes a common generator
whose two integral eigenvalues differ by a unit. L84C2's polynomial
comparison (using that coefficient isomorphism) and the unit-scalar
calculation above put the same difference in the maximal ideal.
This proves the abstract implication as well. On the original
member the internal spectrum proves the separation directly, since
different residual kernels must already differ on a generator.
In particular the result applies for every twisted $i$, every
$i'$ with the same $(j,w)$, and every good row $\sigma$ occurring
in $\mathrm{AtM}$ at both. This supplies the field $c2$ of
$\mathrm{GaloisPreconditions}$ on the same selected member.
% END L84C3

\section{The operator selecting a good cut}

% BEGIN P84C1
\paragraph{P84C1. The polynomial labels of the finite central operator.}
At a generator prime $g_k$ and place $v$, put $A_{kv}=\mathbb F_{g_k}^r$.
The source indexes its local series by
\[
 P(a)=X^r+\sum_{j=0}^{r-1}a_j X^j,\qquad a\in A_{kv}.
\]
Each monic polynomial of degree $r$ occurs once, since its lower
$r$ coefficients recover $a$. This corrects LN1's description
of the indices as ordered root vectors. For instance $X^2+2$ over
$\mathbb F_5$ occurs at $(2,0)$ although it has no root in that
field. Even for split polynomials a root enumeration would introduce
multiplicities absent from this sum.

Let $E_{kv}(a)=e_{\rm Ser,k}(P(a))$ in the complex group algebra
of $G_{kv}=\mathrm{GL}_r(\mathbb F_{g_k})$. These are central
idempotents. The local finite type record sums its series over
their supported monic polynomials to $1$ and gives zero on the
other polynomials. By the bijection above,
$\sum_{a\in A_{kv}}E_{kv}(a)=1$. Embed each factor in the group
algebra of $G=\prod_{k,v}G_{kv}$ by extension at that place.
For a tuple $a$ write $E(a)$ for the product of the embedded
factors. Its coefficient at $g$ is the product of the local
coefficients at the coordinates of $g$. Thus the definition gives
\[
               e_0=\sum_{a:\,\mathrm{goodTuple}(P(a))} E(a),
                \qquad \sum_{a}E(a)=1.
\]
The second equality follows by finite distributivity. Each $E(a)$
is central and idempotent. The trace calculation of LN1 lifts
$|G|E(a)$ through the specified coefficient-ring embedding, using
the primitive $|G|$th root in the given depth-zero record. Summing
over the indicated tuples gives the central integral element $z_i$
with image $|G|e_0$.
% END P84C1

% BEGIN P84C2
\paragraph{P84C2. The selector acts on every type vector above a theta point.}
Assume the four compatibility conditions $\mathrm{Tie4}(D)$,
and take the automorphic constituent, twist and same-point clauses
of the Galois preconditions. Write $Y_i=Z_{\rm good}(i)$.
For every $x\in Z_i$, every automorphic representative $\Pi$
above $x$, and every vector $F$ of its specified type space,
\[
              e_0F=\mathbf 1_{Y_i}(x) F .
\]
First take $F$ in the parameter space of $\sigma$. At each $kv$
the $\mathrm{SerLoc}$ condition says that a local series acting
nontrivially has polynomial equal to $\mathrm{labPar}_{kv}(\sigma)$.
Its lower coefficients form the unique possible local vector $a_{kv}$.
If a tuple differs at any one place, that factor kills $F$, hence so
does its product. Summing all products gives $F$ by P84C1. Thus
$e_0F$ equals $F$ or zero according to the goodness of the
tuple of parameter labels.

It remains to put every $\Pi$ in the indicated parameter space.
At $x\in Z_i$ there is a representative with a base-change list
containing the distinguished twist of the character block
$\mu$ of length $t$. The same-point clause of
$\mathrm{provisionFortyOne}$ identifies the lists of two
representatives of $x$ up to permutation, and $\mathrm{adm1}$
supplies a list for $\Pi$. Move one occurrence of this block
to the end, preserving all multiplicities. Twist the remaining
list by $\eta_T$ and call the result $\Sigma$. Twisting
$\Sigma\mathbin{+\!+}[\mu]$ by $\eta_T^{-1}$ recovers the
permuted list: the twists are inverse on each positive-rank
constituent by $\mathrm{tie0}$ and $\mathrm{adm5}$. If
$x\in Y_i$, start with a good list $\Sigma$ witnessing that
membership. The same permutation argument transfers its
base-change identity to $\Pi$.

For either choice, parameter surjectivity in the
$\mathrm{Pend}$ witness of $\mathrm{Tie4}$ gives a parameter
$\sigma$ with these local values. The defining sum of subspaces
for $\mathrm{MtPar}(\sigma)$ contains the type of $\Pi$,
since $\Pi$ and $\Sigma$ satisfy its two incidence conditions.
The local polynomial dictionary identifies goodness of $\sigma$
with goodness of $\Sigma$ place by place. At $x\in Y_i$ our
chosen list is good. If $x\in Z_i\setminus Y_i$, a good
$\Sigma$ would itself witness $x\in Y_i$. The first paragraph
now proves the assertion for this very $F$.
% END P84C2

% BEGIN P84C3
\paragraph{P84C3. The good cut and the integral commuting operator.}
For every $D,\mathrm{nch}$ satisfying the native
$\mathrm{GaloisPreconditions}(D,\mathrm{nch})$ and
$\mathrm{Tie4}(D)$ there is a good cut, and the uncut datum
satisfies $\mathrm{CommutingEndomorphismCondition}$. Both
conclusions use the same operators. Keep $i$ exactly when
$Y_i$ is nonempty and set
\[
 S_i=\{x:\ \exists\Pi\text{ above }x,\ \exists F\in
       \operatorname{type}(\Pi),\ e_0F\ne0\}.
\]
Each point has a representative with nonzero type space. P84C2
gives $S_i\cap Z_i=Y_i$. With $N_G=|G|$, let
$\alpha_i=z_{O,i}(z_i)$ act on the integral Hecke module, for
$z_i$ of P84C1. Its action commutes with Hecke by the same
$\mathrm{Pend}$ witness. At $x\in Y_i$ its complex scalar on
all type vectors is $N_G$. At $x\notin S_i$ the scalar is zero,
including for $x$ outside $Z_i$, by definition of the support.
The $\mathrm{EigSpan}$ condition passes each scalar to the entire
rational eigenspace at $x$. This proves the operator property
of the cut and gives the commuting endomorphism at every tested
uncut index, since $Z_i\setminus Y_i$ misses $S_i$.

For the last goodness clause, fix $x\in S_i$ and a witnessing
$\Pi,F$. One good tuple has $E(a)F\ne0$. At each place its local
factor acts nontrivially, so one of the irreducible types in that
series occurs in $\Pi$. The goodness conditions give a squarefree
label $P(a)$ with nonzero constant term avoiding the prescribed
unit and its inverse at $j(i)$. For the root exclusions take
any auxiliary prime distinct from $g_k$ in the universal label
test. The $\mathrm{admDZ}$ field supplies zero monodromy and
the inertia polynomial with multiset from $P(a)(X-u)^t$ and its
conjugate-place prescription. A squarefree polynomial over the
perfect residue field is separable. Every repeated root of this
multiset is therefore $u$, and the corresponding unit at the
conjugate place also has order dividing $g_k-1$. The inertia
characters are trivial on wild inertia. Distinct nonzero labels
remain distinct on a tame generator: equality there gives
equality on the dense cyclic subgroup of the tame quotient and
then on inertia by continuity, where the characters distinguish
the labels. Multiplicativity thus transports the order condition
to every repeated inertia eigenvalue. Complete splitting and
the generator-prime size bound give residue norm $g_k$ and
multiplicity threshold two, yielding $\mathrm{Nc}$ as in LN2.

Finally take any other representative above $x$. Its base-change
list agrees up to permutation by the same-point clause. At these
split places the local parameter is given by this list, so it has
the same local condition. All five cut clauses follow, with
$\mathrm{Nc}$ on every representative. T83C3 gives $\mathrm{Tie4}$
on this same cut.
% END P84C3

% BEGIN P84C4
\paragraph{P84C4. The commuting condition on a good cut.}
In any good cut, at every retained index, the cut theta subset
equals the cut good subset. Indeed both memberships restrict
from the uncut point at that index. Goodness gives
$S_i\cap Z_i=Y_i$ and retains all of $Y_i$.
Take $\alpha=N_G\operatorname{id}_{M'}$ on the cut integral module.
This commutes with the restricted Hecke algebra and acts by the
required integer $N_G$ on every good eigenspace. The theta points
at which the commuting condition also asks for vanishing form
the empty set. This proves the commuting component of
$\mathrm{ExtraAlt}$ for every good cut, independently of the
norm-choice facts. It imposes no vanishing condition at the
retained points outside theta.
% END P84C4

\section{Galois preconditions and norm facts on the good cut}

% BEGIN C84F1
\paragraph{C84F1. Residual data and good multiplicities on the cut.}
Fix one datum $D$, norm choice $n$ and good cut $c$. At a retained
localization write $X$ for the old finite set of Hecke points,
$S\subset X$ for the nonempty retained set, $O$ for its coefficient
DVR, and $T\subset O^X$ for the original local Hecke algebra.
The cut algebra is its image in $O^S$. For any $f\in T$ its values
at all points have identical residue. Indeed the evaluation at
each point carries the constants onto $O$. The inverse image of
the maximal ideal of $O$ is therefore a maximal ideal of $T$.
Locality makes these ideals equal, and $f-f(x)$ is in the one at
$x$. It follows coefficient by coefficient that the reduced Hecke
polynomials at $x,y$ agree at every test prime.

A residual rank-$r$ representation is tested by its continuity
condition and its Frobenius characteristic polynomials. A test at
every retained point implies the test at an arbitrary old point:
choose one retained point and use the equality of reduced Hecke
polynomials just proved. Restriction gives the converse. The same
argument with the same residue polynomial witnesses transports
$\mathrm{AtM}$, the condition that a given automorphic label occur
in the residual localization. Thus the old and new good parts in
the complex source are the same subspace.

For $x\in S$ a representative has new point $x$ exactly when its
old point is the underlying $x$. Hence the possible parameter
lists and their complex type spaces are unchanged, as are the
sets $Z$ and $Z_{\rm good}$ after intersection with $S$.
Goodness gives the reverse inclusion $Z_{\rm good}\subset S$.
Consequently the sum of the indicated complex spaces over good
points is the same. T83C2 identifies the full rational eigenspace
at a retained point isomorphically with the original one. Its
dimension, the local multiplicity, is the same; excluded points
make no contribution to the sum weighted by goodness. These
facts transfer each image, kernel and dimension identity of
$\mathrm{ConclusionZ}$, the good-part clause in row C, and its
good-point nonvanishing test. Local uniqueness transfers because
both old $\mathrm{AtM}$ predicates hold whenever their cut versions
hold, and equality of two underlying retained indices implies
their equality as indices of the cut. The residual irreducibility
and tame-parameter tests transfer by the residual and parameter
identifications above.
% END C84F1

% BEGIN C84F2
\paragraph{C84F2. The integral coefficient and row data of the cut.}
Assume the original $\mathrm{GaloisPreconditions}(D,n)$.
A live cut index is live on $D$: its witness is the underlying
point, with the same $Z$ label and local prime. Take at that index
the original smaller coefficient DVR $A$, its embedded Hecke
algebra $T_A$, the integral models of the two representations and
the characters. On $S$ use the image of $T_A$ by restriction and
the restrictions of the same matrices. The $O$-span of their
embedded Hecke algebra is the new $O$-Hecke algebra: applying the
linear restriction map to $O\cdot T_A=T_O$ commutes with sums and
scalar multiplication. Its image is local, and the restriction
of a separator for any two retained points still separates them.
Equations between characteristic polynomials, the chosen changes
of basis and the projected idempotents hold after evaluation at
every point of $S$. For a matrix in the algebra generated by the
representation use its finite expression $\sum f_j r(g_j)$,
restricting each coefficient. In the congruence module the
additional coefficients must vanish at the selected point; they
retain this property after restriction. This transports the
chosen congruence identities integrally, including their original
uniformizer power. These witnesses give the full cut
$\mathrm{Large}$ condition.

The global fields, infinity types, and analytic and geometric
coefficient embeddings of the datum are the same. The
automorphic Galois admissibility test at a retained index uses
the same semisimple realization, with all its local equalities,
its inertia and monodromy comparisons and its de Rham tests.
The signature and temperedness tests therefore restrict at once,
as do the weight and block-character dictionaries, unramified
twisting and depth-zero tests. Combined with C84F1 this gives
the cut admissible datum and its rows A, B, C. The vacuous
auxiliary slot in row A has predicate $\mathrm{True}$. For the
rank test on a size datum on the cut use its identical scale
and numerical fields on the original member. This gives the
original rank bound for precisely that cut input.

Here is a way to transfer the first provision without rerunning
the automorphic construction. In its family clauses a cut
$\mathrm{GoodR}$ can be extended by the given original good
integral realization at every discarded point: its defining
Hecke, continuity and semisimplicity tests are pointwise. All
retained characteristic-polynomial evaluations and parameters
are the old ones. For the clause that characteristic coefficients
belong to the Hecke algebra, take the restriction of the old
coefficient in $T$. For the residual comparisons use C84F1.
For tests conditional on liveness, a live cut index is live
upstairs. For pointwise tests at an unretained automorphic
representative the hypothesis that it has a specified cut point
forces its underlying point to be retained. The family of
rational realizations is obtained by restriction as well.
For the parameter clause $42$ the hypothesis is a nonzero
isotypic row. Its $\theta$-images, its type subspaces and the
definition of a parameter of that row all use the original
global row and automorphic parameter spaces at the underlying
index. They do not use the point-valued option of a
representative. The old positive scalar, theta representative
and parameter list thus serve at the cut, as does cuspidality
on a good row. This completes the first provision.

The cut norm choice is $n_c=\mathrm{nchCut}(c,n)$. Its twisted
place, side, auxiliary set and character ties restrict to the
underlying index, with the old common character witness. The
triple, equal-twist, Fourier-type, inert base-change and
$\lambda$ tests retain their original witnesses. Together with
C84F1 this gives $\mathrm{GaloisPreconditions}(D_c,n_c)$,
where $D_c=\operatorname{cut}(D,c)$. Also $D_c$ has
$\mathrm{CutRec}$: any representative with a cut point has an
old point in $S$, so the nice-local-parameter clause of goodness
applies to that very representative at both generator primes.
% END C84F2

% BEGIN C84F3
\paragraph{C84F3. All labelled de Rham jumps and the admissible jump datum.}
In this and the next three paragraphs work on any $D$ with
$\mathrm{GaloisPreconditions}(D,n)$ and $\mathrm{CutRec}(D)$.
Fix $i,r,x$ with $r$ good and a place $w\mid\ell_i$. Choose an
automorphic representative $P$ over $x$ and a base-change list
$\mathcal P$. At a generator prime, $w$ splits over the real
subfield and its local parameter for $P$ equals that of
$\mathcal P$. The cut hypothesis gives the prescribed tame,
zero-monodromy and inertia-root multiplicity condition $\mathrm{Nc}$ on
this parameter. At any other coefficient prime the parameter
is unramified: at a split place this is the unramified level
condition off the two generator primes and $S_i,Q_i$, and at
an inert place it is in the base-change definition. An
unramified parameter satisfies $\mathrm{Nc}$.

Form the semisimple Galois realization associated to
$\mathcal P$ using automorphic Galois admissibility on its
constituent parts and their prescribed Tate shifts. Its split
Frobenius polynomials agree with those of $r_x$. Density and
semisimplicity give a conjugacy over $\overline{\mathbb Q}_{\ell_i}$.
The admissible comparison for $\mathrm{Nc}$ gives de Rhamness and
the labelled Hodge weights of this realization, and the
conjugacy clause of the comparison transfers both to $r_x$ at
each embedding inducing $w$. Thus its jump multiset is the
prescribed level-zero multiset. Away from the generator primes
the crystalline assertion and the same weight equality are
already in the first provision. Moreover the weights in this
multiset are distinct by regularity of the original infinity
types. If a nonzero irreducible constituent appeared twice
after splitting, de Rhamness of summands and additivity for
direct sums would repeat one of its jumps at any inducing
embedding. This proves distinctness. A common finite splitting
extension for the finite set of points suffices: choose bases of
the absolutely irreducible summands of their semisimple
representations over the algebraic closure and adjoin all
entries of the change-of-basis matrices. This proves the two
native $g2$ assertions.

For $\mathrm{admJ}$ take an arbitrary size datum $L$ with
$\mathrm{AdmPar}(L)$. For a de Rham representation the member's
Weil--Deligne comparison identifies crystallinity with unramified
Weil--Deligne parameter. Conversely a crystalline representation
is de Rham by the crystalline/de Rham period inclusion. These
comparisons may be tested over a finite coefficient extension:
the inverse images of $\mathrm{GL}_a$ of the finite extensions
cover the compact Galois group by countably many closed subgroups.
One is open by Baire, and adjoining the matrix entries of its
finitely many coset representatives suffices. Entrywise
restriction to that finite field preserves continuity.
For the weight dictionary distinguish the chosen archimedean
embedding over the real place and its conjugate. The CM type
contains exactly one of the two. Use the member's weight
parameter at that embedding and its block character there; the
common finite-order factors do not change the labelled jumps of
the character. The arithmetic character of the Galois realization
and the chosen automorphic character agree at the split
Frobenius tests and hence on the Galois group by density.
This gives the two signs and the common shift of the native
dictionary (the global orientation parameter is false). More
explicitly, with $\sigma=\operatorname{sgnOf}(\mathrm{true},b)$
for the embedding choice $b$, the level-zero multiset is
$\{\sigma u+c_\tau:u\in \mathrm{PP}(D,i,\rho)\}$ with
multiplicities, and the first character has its single jump at
$-\sigma e_\rho(\mathrm{liftIdx}(i))-(t-1)/2+c_\tau$.
In the local reading, compose the given embedding of the
completion with $E\to E_w$. It induces $w$. Under a rank-one
twist with labelled weight $c$ the degrees shift by that
weight, and duality reverses the grading. With the prescribed
Tate twist the reflected label is $c-1-a$ at degree $a$.
Additivity on a given direct-sum decomposition and the test
$\mathrm{gr}^a\ne0$ if and only if its dimension is nonzero
give the ordinary and dual summand tests over every admitted
coefficient extension. These are the $\mathrm{AdmRead}$ clauses;
the two $g2$ conclusions above and the given $\mathrm{AdmPar}$
supply the remaining entries of $\mathrm{AdmJAlt}(L)$.
% END C84F3

% BEGIN C84F4
\paragraph{C84F4. The tame polynomial and vanishing of the local cohomology in G5.}
Use local finiteness, local Tate duality and the local Euler
characteristic theorem. At an index satisfying $\mathrm{HG}$ choose
a good point and its good parameter. Row C supplies a residual
rank-$r$ realization there. Its Frobenius tests at that point
are those at every point by C84F1's local Hecke calculation.
We may thus choose a global residual part $\bar\rho$. For any
such part, every good family $r$, and every $w\in S_i\cup Q_i$,
the base-change comparison at $w$ identifies inertia polynomials
and gives zero monodromy. Indeed $w\nmid\ell_i$ and splits.
Use the same semisimple conjugacy as in C84F3. Its
Weil--Deligne relation at $w$ gives the full dimension of the
local parameter. At a generator place the needed $\mathrm{Nc}$
comes from the cut; at $Q_i$ and the specified auxiliary primes
the given monodromy/inertia type and its full dimension give
$\mathrm{Nc}$. In the split case of the G5 statement that type
has $r$ trivial characters and $t$ copies of $\chi_1$. Vanishing
monodromy gives a finite inertial image; the simultaneous
character splitting is over the fraction field because both
characters are defined there. Consequently
$(r_x(\gamma')-1)(r_x(\gamma)-\chi_1(\gamma))=0$ for every
two inertia elements there.

Set $z=\overline{\chi_1(\gamma_0)}$ for the prescribed tame
generator at $w\in S_i$. At either generator prime $z$ is excluded
from the roots of $\bar\rho(\gamma_0)$ by residual admissibility:
$w\nmid\ell_i$ supplies its unequal-prime guard. At the remaining
auxiliary prime, in the split calculation above the full residual
polynomial is $(X-1)^r(X-z)^t$. The global residual
factorization is $\operatorname{char}_{\bar\rho}(X)(X-z)^t$
because the cyclotomic character is trivial on this inertia.
Cancel the monic factor. Here the auxiliary inertia character
has image of order $15$ and the designated generator has
primitive value. Indeed wild inertia at this prime, whose
residue characteristic is $1$ modulo $15$, maps trivially to
the cyclic image of order $15$. A primitive fifteenth root has
unit difference from one since $\Phi_{15}(1)=1$. Thus in either
case $z$ is a nonroot. For every integer $e$ the residual twist
$M=\bar\rho\otimes\bar\chi_1^{-1}\bar\epsilon^{e}$ has
$\gamma_0-1$ invertible. This kills its invariants. The same is
true on its Tate dual: on inertia the cyclotomic character is
trivial and inverse transpose preserves the absence of eigenvalue
one. Local Tate duality kills $H^2$; the finite local Euler
characteristic away from $\ell_i$ gives
$|H^1|=|H^0||H^2|=1$. These are all the required twists.

At a generator place in $S_i$ form the polynomial of
$r_x(\gamma_0)$ simultaneously in $T[X]$ by the coefficient
clause of the first provision. Its reduction has root $z$ of
multiplicity exactly $t$ by the nonroot calculation. Its
characteristic-zero multiplicity at $\chi_1(\gamma_0)$ is at
most $t$, since reduction can only increase this multiplicity.
Apply the generator-prime clause of the input
$\mathrm{MonodromyInertiaType}$ at $w\in S_i$ over $g(k)$:
its two tests are precisely that $\gamma_0$ generates tame
inertia and this characteristic-zero multiplicity is at most
$t$. It supplies a decomposition of the local parameter into
parts of dimensions $r,t$. The second part has zero monodromy
and characteristic polynomial $(X-\chi_1(\gamma))^t$ for
every inertia element $\gamma$. The Weil--Deligne comparison
for a direct sum gives complementary inertia-stable spaces.
On the second one finite inertia acts semisimply in
characteristic zero, hence every $\gamma$ acts as the scalar
$\chi_1(\gamma)$. Divide the universal characteristic polynomial by the
monic $(X-\chi_1(\gamma_0))^t$ in $T[X]$. Its remainder is zero
at all $x$, hence zero in $T[X]$. The quotient $P$ is the
characteristic polynomial of the complementary block at each
$x$. Cayley--Hamilton there, and the character action on the
other block, give
$P(r_x(\gamma_0))(r_x(\gamma)-\chi_1(\gamma))=0$ integrally.
Its value at $\chi_1(\gamma_0)$ is a unit in the local algebra
by the residual multiplicity. This supplies one polynomial for
all points and all inertia elements.

Finally apply the idempotent-choice functional of G5 to the
corresponding matrix relation on its congruence generator
$y_\gamma$, a right multiple of $r(\gamma)-\chi_1(\gamma)$
with the prescribed scalar normalization. The left operator
$P(r(\gamma_0))$ belongs to the indicated Hecke span. Its
character acts on this functional by the displayed unit, so
its vanishing on the product forces vanishing on $y_\gamma$.
In the auxiliary case use $r(\gamma')-1$ in place of the
polynomial. Choose $\gamma'$ giving a primitive fifteenth root;
its character minus one is a unit and the two-element split
identity above gives the zero product. This proves all clauses
of $\mathrm{ConclusionG5}$, including the integral functional
test.
% END C84F4

% BEGIN C84F5
\paragraph{C84F5. A character eigenline detects the full rank-one string in G6.}
Take a good index, a good integral family $r$, and a point $x$ at
which $r_x$ has a nonzero $\chi_1$ eigenvector over the fraction
field. Choose a representative and any one of its base-change
lists $(a_k,\pi_k,d_k)_k$. For the $k$th part take the admissible
centered Galois realization $R_k$ of rank $a_k$. For the required
algebraic twist, the top weight of this string at infinity is the
weight of $\pi_k$ translated by $(d_k-1)/2$. Its multiset is
contained in the regular algebraic weight of the entire list, so
it has no repeats and lies in the required translated integer
lattice. This is the parity and regularity condition for taking
the centered twist of the attached cuspidal realization in the
given automorphic Galois admissibility hypothesis. The polynomial
of $r_x(g)$ is that of the direct sum
\[
               \bigoplus_k\bigoplus_{0\le j<d_k}
                                     R_k\epsilon^{-j}.
\]
Indeed both polynomials at every split unramified Frobenius away
from the finitely many excluded places are the factors of the
same base-change parameter, with $\epsilon(\mathrm{Frob}_w)=q_w$.
All coefficients are continuous, and these Frobenius conjugacy
classes are dense. Both sides are semisimple over the algebraic
closure. Brauer--Nesbitt and projection of the nonzero eigenline
therefore give a $\chi_1$-eigenvector in at least one shifted
summand, say $(k,j)$.

At a split place at which $\pi_k$ and the distinguished character
are unramified, the vector gives eigenvalue
$q_w^{j}\iota(\chi_1(\mathrm{Frob}_w))$ in $R_k$. Compare its
absolute value with the normalization of the centered parameter.
For the corresponding unshifted cuspidal Satake root $\alpha$,
this gives $2\bigl|\log_{q_w}\lvert\alpha\rvert\bigr|=|2j+t-d_k|$. The native open cuspidal bound confines its absolute
value strictly between $q_w^{-1/2}$ and $q_w^{1/2}$.
The bound applies also when the integral parity condition for the
cuspidal part requires the fixed half-weight character: use the
admissible twist and remove its known norm in the Satake and
infinity parameters. Since $2j+t-d_k$ is an integer and $q_w>1$,
we obtain $d_k=t+2j$.

For each $k$ let $A_k$ be the multiset of dimensions in the
residual semisimplification of a stable lattice of $R_k$, after
extending coefficients to a common algebraic closure of the
residue field. Such lattices exist: the matrices descend to a
finite local coefficient field as in C84F3, and the span of the
translates of a lattice lies in a bounded multiple of it by
compactness. Each $A_k$ has positive entries of sum $a_k$.
Reduction of the polynomial identity gives
\[
            \sum_k d_k A_k=\{r\}+t\{1\}.
\]
In fact the comparison on the right is the residual polynomial
of C84F4's absolutely irreducible rank-$r$ part and the $t$
characters. Semisimplify after reduction on both sides; twisting
preserves irreducibility and dimension. The polynomial at every
group element determines the semisimple constituents: the
residual lists factor through a common finite quotient, and on
its prime-to-$\ell$ elements the matching eigenvalue multisets
give the same Brauer character. Taking dimensions with
multiplicities gives the stated equality.

Here $r>1$ and $t\ge2$. Our selected $k$ has $d_k\ge t$.
An entry larger than one in $A_k$ would occur at least twice
on the left and at most once on the right. Every entry is
therefore one, and there are at most $t$ available copies,
forcing $A_k=\{1\}$, $a_k=1$, $d_k=t$ and $j=0$.
The eigenline is now the whole $R_k=\chi_1$. The prescribed
one-dimensional automorphic character $\mu_i\eta_T^{-1}$
has the same normalized realization: its Frobenius values
agree with $\chi_1$ by the given character tie, hence agree
on the whole group by density. Undoing the common local twist
identifies the corresponding local automorphic characters almost
everywhere; admissible multiplicity one identifies them globally.
This identifies the required part in the base-change list, so
$x\in Z_N$.

Finally every parameter $\Sigma$ at $x$ gives another such list
by appending the prescribed rank-one part and twisting. Apply
the residual multiset calculation to this list. Its appended
part uses all $t$ copies of $1$, leaving the singleton $\{r\}$.
Every remaining part has nonempty $A_k$ and $d_k\ge1$. Exactly
one remains, with $d_k=1$ and $a_k=r$. Undoing the twist gives
$\Sigma=[(r,\pi,1)]$. This proves both native conclusions for
every tested eigenline, including the universal parameter
quantifier.
% END C84F5

% BEGIN C84F6
\paragraph{C84F6. Coefficient duties and the actual cut norm-choice facts.}
The four coefficient duties use the same residual representation,
ambient point family and smaller coefficient model as above.
The first is $\mathrm{OLarge}$ from $\mathrm{Large}$. For distinct
constituents, use C84F3 on a finite splitting field. The
$\mathrm{OSplit}$ clause of $\mathrm{Large}$ splits the original
semisimple good representation into absolutely irreducible pieces
over the original fraction field. If two disjoint pieces there
were isomorphic, they would remain so after scalar extension,
contradicting C84F3. This supplies the required condition over
the original coefficients.

For the design test the given residual matrix span on the
rank-$r$ summand is the full matrix algebra; after any residue
extension a stable subspace for all these matrices is either
zero or the whole space, by the elementary matrices. This is
absolute irreducibility. At a type-prime test the given rank
decomposition instead identifies the residual part by dividing
the Hecke characteristic polynomial by the prescribed character
factors at split Frobenius places. C84F1's residual comparison
makes this a residual part at every point. The original
residual irreducibility test now supplies the same conclusion.
In either case the chosen point and the given decomposition
witness liveness, so $\mathrm{Large}$ supplies its smaller-ring
model including the projected block $e_B$, the complementary
rank-one matrix $e_1$, the regular element $X_0$ and the
specified $g_0$ in inertia.

Here are the integral issues in using that model for the
representation in the test. Extend the smaller ring into $O$;
the chosen block polynomial, monic of degree $t$, has $t$
distinct roots and already splits over the specified fraction
field. Its roots lie in $O$ by integrality. Over each point take
the $t!$ orderings of those roots. Over $T$ the successive monic
root adjunctions have bases of lengths $t,t-1,\ldots,1$.
Their product basis is independent in the ring of functions on
these ordered tuples: over each fraction-field point the roots
are distinct, giving the product of fields by successive
interpolation, and the original $T$ embeds in their product.
Thus the $T$-algebra generated by the root functions has its
required integral free basis, by the same monic-division argument
as D79G5, with distinctness needed only over the fraction field. At the distinguished lift order
them as $\chi_c(g_0)$, $1\le c\le t$. Multiplying the other
linear factors of $X_0$ in its $e_B$ block selects the first
line multiplied by the product of its root differences.
At every ordering this product of factors has rank one and its
square is its scalar root-product times itself, by evaluating on
each distinct eigenspace in the block. The characteristic and
block determinant coefficients belong to the Hecke span by the
word expansion of D79G2. These identities hold integrally by
injectivity into the fraction field. The residual
absolute irreducibility gives all integral rank-block matrices
as evaluations of the chosen algebra: reduction spans the
matrix algebra by the density theorem, and Nakayama lifts the
span. The resulting ordered-root record has the original
$g_0$ and the prescribed product in its length clause.

For a different tested good family, density gives identical
rational characteristic polynomials pointwise, hence conjugate
semisimple representations over the fraction field. Cancelling
the prescribed character summands gives the conjugacy of the
rank blocks at the tested point. Transport every matrix in the
generated algebra by evaluating its original Hecke-linear
group-algebra expression in the new integral matrices. This
defines both it and the inverse transport integrally because
conjugation over the fraction field detects equality of two
expressions. The rank-block action on its integral lattice
still gives every matrix: the conjugacy $Q$ on the rank block
carries every old elementary matrix into an integral matrix
(the conjugated expression acts integrally by a sum of the
new rank representation matrices). Thus $Q M_r(O)Q^{-1}$
lies in $M_r(O)$. Conjugation on the full matrix space has
determinant one; both orders are lattices of the same covolume,
so containment is equality. Character evaluations, ranks,
determinants and the ordered-root functions stay the same.
The element $g_0$ thereby retains its required wild cyclotomic
condition at a type prime. These are respectively the
ordinary and type-prime coefficient duties.

Apply C84F3--C84F6 to $(D_c,n_c)$ of C84F2. In the norm-choice
record, the second provision consists precisely of C84F4,
C84F5 and the good-point identities of C84F1. Insert the $g2$
and parameter conclusions of C84F3 and the coefficient duties
just proved. Every remaining entry is the corresponding cut
precondition of C84F2. This gives
$\mathrm{NormChoiceFacts}(D_c,n_c)$ for every good $c$ on a
member with the stated original preconditions. The arithmetic
principles used here are the given admissibility and row
hypotheses, Chebotarev, and local finiteness, Tate duality and
the local Euler characteristic formula. In particular D79G7
used G5 and G6 in its proof on a supplied hub. The two
computations above supply them in the reverse direction
required here from the preconditions and $\mathrm{CutRec}$.
% END C84F6

\section{The denominator with the norm-choice family hypotheses}

% BEGIN V84A1
\paragraph{V84A1. The information retained by the alternate hypotheses.}
Fix a member $D$, its norm choice $n$, and its
$\mathrm{NormChoiceFacts}(D,n)$. At every localization $i$ and for
every good integral family $R$, the following data are available.
\begin{enumerate}
\item The coefficients of $\det(1-zR_x(g))$ belong to the Hecke
algebra as functions of $x$, the residual characteristic polynomial
at $g$ is independent of $x$, and $R$ is unramified outside the
places above $p=\ell_i$ and the two sets $S_i,Q_i$. After one finite
coefficient extension, every irreducible subrepresentation at every
$x$ is absolutely irreducible and disjoint irreducible
subrepresentations at the same $x$ are inequivalent.
\item Off the two type primes, every $R_x$ is crystalline above $p$,
with the specified level-zero jumps. At either type prime every
$R_x$ is de Rham there with the same labelled jump formula.
At parameters satisfying $\mathrm{AdmPar}$ one has the sign
dictionary, the four multiplicity readings and the
crystalline--Weil--Deligne tie of $\mathrm{AdmJAlt}$.
\item The quadratic inertia identities and the Hecke polynomial at
a type place in $S_i$, the off-support criterion for a global
$\chi_1$-line, and the good-row image and rank identities have
the original $\mathrm{Conclusion2}$ quantifiers.
\end{enumerate}
Indeed $\mathrm{provisionFortyOne}$ supplies $G1,G2a,G3,G4$ in the
alternate definitions. Take the second factor of $G3$ for the
residual polynomial. Its first factor requires a guard, but the
second factor of $\mathrm{DeRhamJumpsDistinctCondition}$ gives
the entire characteristic-zero distinctness assertion for the
chosen good family without that guard. The first factor of the
latter condition gives exactly the type-prime reading. $G2a$ is
the off-type crystalline and jump assertion itself. Take the
remaining reading fields from $\mathrm{admJ}(L)$ at the supplied
parameter. For (3) project to $\mathrm{ConclusionG5}$,
$\mathrm{ConclusionG6}$ and $\mathrm{ConclusionZ}$, respectively.

For clarity about the comparison, the seven fields extending the
core to norm-choice facts (localization uniqueness, large-prime
conditions, coefficient duties, norm-choice tie, good triples, same
twist and goodness of all the specified signatures) are unchanged.
In the core, $\mathrm{Conclusion2}$, residual irreducibility of the
rank block, weight and character dictionaries, unramifiedness of
the twist, depth-zero condition, spherical base change and
$\mathrm{LamUsed}$ are unchanged. The changed entries are the
row record, admissibility, first conclusion, residual tame labels,
de Rham conditions and $\mathrm{admJ}$. Inside the row record
the changes are in its first part, concerning Galois attachment,
the cuspidal bound and the former all-place invariant condition.
The alternate attachment gives a relation of local parameters
away from $p$, giving equality whenever the automorphic parameter
is unramified. Its clause above $p$ has a rank-one or local-parameter
guard. The alternate cuspidal bound gives temperedness in rank one
and $q^{-1/2}<|\alpha|<q^{1/2}$ on unramified roots. The former
all-place invariant condition is replaced by the true predicate.
The alternate residual tame label assertion has a squarefreeness
premise at all the type labels. None of (1)--(3) requires removing
that premise. Finally $\mathrm{ExtraAlt}$ keeps the same
commuting-endomorphism field. Its paired-character inequality
has the same right side and omits the argument $\mathrm{SetSQ}$.
It therefore implies the ordinary paired inequality by applying
it to the other arguments. These projections specify the inputs
used below without an implication between the complete core records.
% END V84A1

% BEGIN V84A2
\paragraph{V84A2. Constructing the alternate setup on its own member.}
Assume $\mathrm{Hypotheses}_2$, $\mathrm{TwGood}(\mathrm{Fm})$ and
$\mathrm{FactsFAlt}(\mathrm{Fm},C,C_T)$. For each real $C_2$ there is
a cutoff such that for every subsequent $x>\mathrm{Fm}.x_1$, each
member at $x$ satisfying $\mathrm{ShaFinOdd}$, and every prescribed
finite set $A$, there are $H,F,\mathrm{Cmp}$ and a hub $X$ on this
member with
\[
 X.\mathrm{exc}=A,\qquad \mathrm{AdmJAlt}(L),\qquad
 \mathrm{SetupAlt}(L,n,H,X,C_2).
\]

To construct the point representations, fix a regular algebraic
conjugate-self-dual cuspidal part of rank $a\ge1$. The alternate
attachment supplies an actual continuous semisimple representation.
Discard the finitely many places over its coefficient prime and
the finite set of ramified automorphic places. At each remaining
place its normalized parameter is unramified because the
normalizing twist is unramified. The equality clause of the
local relation now identifies this parameter with the attached
Weil--Deligne parameter. It too is unramified. By the tied
Frobenius polynomial, the trace of this representation at an
arithmetic Frobenius equals the sum of the Satake parameters
times $(\mathrm Nw)^{(1-a)/2}$ under the specified coefficient
identification. This proves precisely the cofinite trace input
DFD4 supplies in the ordinary construction.

Choose the good representation family promised at each index by
the first existence clause of $\mathrm{provisionFortyOne}$.
Using these trace representations, the two parity cases of DFD5
give all shifted blocks. DFD3 and DFD6 identify their traces on
the retained Hecke point and produce its character splitting.
The rational descent and saturation in DFD1--DFD2 give a
continuous complement over this very integral coefficient ring.
DFD7--DFD8 assemble the point family and its canonical hub for
the given twist, using zero as error function and $A$ as the
constant exceptional-set function. This is the construction of
DFD9 with the trace input proved above.

Use the three cutoffs of F78D2 for the family parameter, small
primes at constant one and the good twist. Choose $H,F,\mathrm{Cmp}$
by L78W13 before forming that hub. The good twist gives its twist
index. The first conclusion about theta lifts and its good-row
image, used by DFD10, are respectively $\mathrm{Conclusion42}$
in the alternate first conclusion and $\mathrm{ConclusionZ}$ in
the second. For the actual point representations the norm-choice
tie, good triples, same twist and residual irreducibility are the
fields in V84A1. Thus SSC15--SSC16 and DFD11 give killing on
$X.\rho$. The family size cutoff gives $8\le L.\mathrm{sc}.L$ and
$60\le D.\mathrm{host.splitBound}$. DFD13 gives $\mathrm{AdmPar}$
and the polynomial parameter bounds, hence $\mathrm{AdmJAlt}$ by
the member's norm-choice facts. These data and the supplied
finiteness fill every field of $\mathrm{SetupAlt}$.
% END V84A2

% BEGIN V84A3
\paragraph{V84A3. The jump, global Hecke and type-prime tests.}
Fix any alternate setup, its actual $R=X.r$, $\rho=X.\rho$ and
the native member parameters. The following verifications also
hold for arbitrary good $R$ and specified rank-$r$ complements
when these are quantified in the corresponding test.

First, S80J1--S80J2 give all seven jump conditions. Their
calculation uses the label and the sign in
$\mathrm{AdmDict}$ and $\mathrm{AdmRead}$ of V84A1. At a non-type
prime, the crystalline comparison gives the required de Rham
reading. At a type prime, use the supplied de Rham clause.
Subtract the same one-dimensional character weight, then subtract
$\{0,\ldots,t-1\}$ at the split rank block. In both orientations
the resulting rank list and its dual $z\mapsto-1-z$ omit zero.
At a non-type prime, the full family is crystalline, and the
character and rank block are crystalline direct summands at the
chosen split point. The same holds after twisting and coefficient
extension by S80J3. Any rational invariant in the rank block or
its Tate dual would supply a trivial de Rham subobject and hence
a jump of weight zero by strictness. This proves both required
local vanishings above $p$ on this data.

At an index with good support the Hecke coefficient and residual
polynomial conditions of V84A1 and the identical $\mathrm{rhoBar}$
field give the full integral matrix span of D79G1. The coefficient
duties supply the stated decomposition over the coefficient field
in S80J3. At a point outside $Z_i$, a global fixed vector after
dividing by $\chi_1$ would supply the $\chi_1$-line tested by
$\mathrm{ConclusionG6}$. Outside $S_i,Q_i,p$, $G4$ gives the
unramified action, including on the character summand at the
chosen point. In $Q_i$ and at the auxiliary prime in $S_i$,
$\mathrm{ConclusionG5}$ gives the quadratic inertia identity;
at a type place in $S_i$ it gives the annihilating Hecke
polynomial and its unit value. The tame nonidentity and auxiliary
order-$15$ unit calculation of D79G7 depend only on the specified
characters. Thus D79G1--D79G7 give the same determinant, local
and off-support inputs in their respective ranges.

For the type-prime estimate assume also the paired-character
part of $\mathrm{ExtraAlt}$, $0\le C$, $C_T\le b$ and the late
family size bounds. At $p=g_k$, take $s=\lfloor3b\rfloor+2$ and
$S=0$ as in T79W1. The log-ratio and size inequalities in Y79T1
and T79W1 depend only on the two type primes and the member
sizes. At the completely split prime, T79W2 gives the opposite
signs at conjugate places and the sets of jumps for every $R_x$
by the type-prime reading just proved. These are
$\mathrm{SignOK}$ and $\mathrm{HostJ}$. Every $\pi\in Z_i$ has
its canonical semisimple splitting. Killing gives the finite
Selmer predicate at every good point and $\mathrm{ShaFinOdd}$
gives finiteness at $0<e<t$. Substitution in the alternate paired
inequality yields exactly the bound for $\eta_{\rm loc}$ in Y79T1.
In the empty good-support case both the good eigenspace index
and the sum vanish as in that proof. Y79T2--Y79T4 then give
$\mathrm{StarH}$ with the original multiplicities and type
constant $2000(C+1)+5C_1$, where $C_1$ is the uniform parameter
bound. Here the multiplicity and rank bounds used in the summation
follow from $\mathrm{Conclusion42}$, $\mathrm{ConclusionZ}$,
$\mathrm{adm8}$, localization uniqueness and the rank field by LN6
and J78H1. The discarded all-place condition does not occur in
this substitution.
% END V84A3

% BEGIN V84A4
\paragraph{V84A4. The small-prime family bound for every alternate setup.}
Under $\mathrm{Hypotheses}_2$, every family satisfies
$\mathrm{BoundAtTwoThreeFamilyFonAlt}$. In fact a single pair
$(C_s,x_0)$ bounds all its good-point rows at $p\mid Q_S$ or
$p=3$ whenever the quantified alternate setup is supplied.

Let $C_1\ge1$ be the parameter constant of A78P3 for this family.
Take
\[
 C_s\ge18\cdot6000 C_1+2\cdot30000(3C_1+9)+5C_1+1 .
\]
Fix the cutoff for these parameter and size inequalities and
$\log x\ge1$ before any member, functor, hub or $C_2$.
Now take arbitrary $H,F,\mathrm{Cmp},X,C_2$ and a witness
of $\mathrm{SetupAlt}$. At a twisted index with a good point the
setup gives a genuine point family, its killing, the norm-choice
facts, the splitting bound and the finiteness at $0<e<t$ for
every prime, including $2$. Specializing the norm-choice facts
at its admissible parameters supplies the reading data of V84A3.
The argument for this assertion uses neither
$\mathrm{ExtraAlt}$ nor a good twist of the whole family.

For these primes G82F1 gives the split declaration, the levels
$s=\max(1,v_p(Q_S))$ and the cyclotomic inertia element. It uses
only the host and splitting bound. Each prime is outside the two
type primes. V84A3 therefore supplies the full ambient and dual
crystalline families on each tested coefficient extension, their
simple labelled jumps and the rank-block invariant vanishing.
In G82F2 construct the integral matrix span, projectors and
auxiliary inertia relations from the Hecke and off-support tests
of V84A3. Their root differences at the chosen lift have valuation
$\sum_{1\le j<t}(s+v_p(j))$. The small-prime calculation uses
the declaration just chosen and the actual $R,\rho$ in this
setup. The row constant $C_2$, hub exceptional set and stored
middle-range lengths do not enter it.

More precisely, L82K5 supplies the integral lifts for all
surjections onto the rank and character columns and for the
dual row maps, since the crystalline subquotients have the
inherited height intervals. T82K6 supplies the torsion and
invariant estimates on those same representations. One may
use the compatible comparison of L78W13 on $D$ for these local
calculations: the maps and finite local groups in G82F2 are
defined on the Galois lattices of $X$, independently of the
presentation used to calculate them. The determinant argument
and descent in G82F2 now give $\mathrm{Bd}$ with its displayed
error $\mathcal E$. By G82F3,
\[
 \mathcal E\le6000 f m h_B s+
    30000 f^2t^2 n_V\bigl(s+\log_p(2h_B)+1\bigr).
\]
G82F4 bounds this by $\mathrm{midF}(L,C_s)$ with precisely
the choice above. At $p=2$ the level is at least two because
$4\mid Q_S$; at $p=3$ outside $Q_S$ the setup's splitting bound
suffices. These include exactly the two alternatives in
$\mathrm{BoundAtTwoThree}$. Keeping also the other primes of
$Q_S$ yields $\mathrm{StarF}$ on this $X.\rho$ by G82F4's
auxiliary table of zero place lengths and non-place entry
$\mathrm{midF}(L,C_s)$.
% END V84A4

% BEGIN V84A5
\paragraph{V84A5. Fifteen simultaneous fields and the alternate-family denominator.}
Under $\mathrm{Hypotheses}_2$, suppose the given family has
$\mathrm{TwGood}$ and $\mathrm{FactsFAlt}(\mathrm{Fm},C,C_T)$
with $C\ge0$. Then it satisfies $\mathrm{ConclF}$ and
$\mathrm{LinkN}$ with their original late-member quantifiers.

Set $C_2=122$. Choose the common table constant as in A83G2,
including the lower bounds in V84A3--V84A4 and $C_T$. At this
constant choose the cutoff meeting all the size bounds
(including $C_T\le b$), the local and global row bounds and the
cutoff of V84A2. At any subsequent member satisfying
$\mathrm{ShaFinOdd}$, prescribe its arithmetic exceptional set in
V84A2. Keep that very hub in the following verifications.
\begin{itemize}
\item The canonical table, exceptional-set equality, admissible and
polynomial parameters give $\mathrm{num},\mathrm{exc},\mathrm{par}$
by F78D3, X78A3 and A78P3. Positivity is the choice of the
constant and $\mathrm{small}$ is DFD14 at this constant.
$\mathrm{pass}$ follows from LN5 with the given commuting
endomorphisms. Multiplicity and rank were checked in V84A3 and
killing is part of V84A2.
\item For $\mathrm{bb}$, the hypotheses on the twisted rank
lattice at the declared columns in S79P1--S79P5 are its
crystallinity, jump bounds and absence of zero from V84A3.
J79H1's exponent $120$ thus gives the coarse bound $7200$ on
this canonical hub by J79H2. The type field is V84A3.
\item At a large-prime test, fix any of its good families,
points and complements. The full span, auxiliary inertia and
off-support conditions of V84A3 give D79G8's global inputs.
The guard excludes the type primes. The jump lists, crystalline
condition and residual and determinant data in V84A3 are exactly
the member inputs to S80J3--S80J6. J80H1--J80H2 therefore give
the three local cohomology assertions and the congruence estimate
for this arbitrary complement. With the unchanged weight dictionary
and the hub's exceptional set, G79E2 gives $\mathrm{large}$.
\item The primes $q,d$ and the guarded middle range are also
outside the type primes. For $q,d$ apply J81H2 on the rank
and character cuts formed with V84A3's global data. The rank
lifts and character bounds there depend on the crystalline
jump lists and the local fields; both are those already verified.
For a middle-range prime use the same Hecke and inertia data for
the common columns of G82C5 or the empty columns of G82I1.
At its places the cut lattice, labelled weights and finite
cohomology are supplied by the canonical table and V84A3.
The split estimate G82C5 and inert estimates B82R3 and A83G1
then give the actual $\mathrm{StarAD}$ on this hub. Their
constants are the prescribed row constants at $C_2=122$.
This gives $\mathrm{ram}$ and $\mathrm{mid}$. The small-prime
table of V84A4 gives $\mathrm{two}$ for the very same $X.\rho$.
\end{itemize}
This constructs every field of $\mathrm{Concl}$ for the chosen
$X$ within the uniform cutoff, hence $\mathrm{ConclF}$.
Apply the table export of A83G3 to these fields. For the same
member and each allowed lift, under summability of its
$\mathrm{vvTerm}$, the weighted $\eta$-sum satisfies
$\mathrm{LinkN\_at}$ with constant $2000(|C_s|+1)^3$.
Both constants and any enlargement of the export cutoff precede
the member. This proves $\mathrm{LinkN}$ on the supplied alternate
family, including when that family is the one obtained by cutting.
% END V84A5

\section{The congruence ideal at a type prime}

% BEGIN P84D1
\paragraph{P84D1. The ideal on the retained points.}
Write $O$ for the coefficient discrete valuation ring, $p$ for
the residue characteristic, $e=\ell_O(O/pO)$, and $X$ for
the original finite set of Hecke points at the chosen index.
On a good cut let $S\subset X$ be the retained set and put
$Y=S\cap Z=Z_{\rm good}$, $B=S\setminus Y$.
For $U\subset X$ let $T_U$ denote the image of the original
Hecke algebra $T$ in $O^U$. For every $\pi\in Y$ put
\[
 J_\pi=\{h(\pi):h\in T,\ h(b)=0\text{ for every }b\in B\},
 \qquad c_\pi=\ell_O(O/J_\pi)/e .
\]
This set of values is an ideal since $T$ contains the scalars.
Restriction onto $T_S$ is surjective, and vanishing off $Y$ in
$S$ means exactly vanishing on $B$. Thus $J_\pi$ equals the
native evaluation ideal on the cut. In $X$ the complement of $B$
is $Z\cup(X\setminus S)$, the native cut locus, so the same
$c_\pi$ also equals the original $\mathrm{vccCut}(i,S,\pi)$.

Apply the coordinate filtration inequality LN4 to $Y$ on the
cut lattice. P84C4 identifies the good and theta subsets there.
Equivalently the general endomorphism index bound applies with
the identity and scalar $1$, whose valuation is zero. With $m_\pi$
the cut eigenspace multiplicity this gives
\[
                  \eta'_{\rm loc}(i)\le\sum_{\pi\in Y}m_\pi c_\pi .
\]
This use of the identity is in the general length inequality.
P84C4 separately supplies the prescribed group-cardinality
scalar for the commuting condition in the family record.
% END P84D1

% BEGIN P84D2
\paragraph{P84D2. An exact calculation from two integral restriction images.}
Write $d(U)=\ell_O(O^U/T_U)$, with $d(\varnothing)=0$.
The separation condition on $T$ makes every such length finite.
Indeed for each $x\in U$ and each $y\ne x$ in $U$ choose a
Hecke function taking distinct values at $x,y$. Subtract its
value at $y$ and multiply these functions over $y$. The result
has support $x$ on $U$ and nonzero value at $x$. For
$U=\{x\}$ the empty product is $1$. These elements span the
coordinate lines over the fraction field, and $O$ is Noetherian.
For every $\pi\in Y$ we have
\[
           e c_\pi=d(B\cup\{\pi\})-d(B).
\]
In fact the map $O\to O^{B\cup\{\pi\}}/T_{B\cup\{\pi\}}$
supported at $\pi$ has kernel $J_\pi$. Its image is the kernel
of restriction to $O^B/T_B$: a representative in that kernel
agrees on $B$ with some original Hecke function, whose subtraction
leaves just its coordinate at $\pi$. Restriction is surjective
by extension of functions. Additivity of lengths proves the
formula. Each $d(U)$ is the sum of the valuations of the
Smith invariant factors of a matrix of integral module generators
of $T_U$, with the uniformizer of $O$ assigned value $1$.

There is also a bound needing only pairs. Let $a_{\pi b}$ be
the least valuation of a nonzero $h(\pi)-h(b)$ for $h\in T$.
It exists by separation and discreteness. Multiplying one
minimizing difference function $h-h(b)$ for each $b\in B$
produces a function whose value at $\pi$ has valuation
$\sum_{b\in B}a_{\pi b}$. Therefore
$e c_\pi\le\sum_{b\in B}a_{\pi b}$. In particular the
congruence lengths vanish if $B$ is empty, or, more generally,
whenever each $b\in B$ can be distinguished from $\pi$
modulo the maximal ideal. On a general good cut the set $B$
may contain points, and the exact difference above tests their
contribution.
% END P84D2

% BEGIN P84D3
\paragraph{P84D3. The point bound implies the paired condition with the same constant.}
Consider any of the tests in
$\mathrm{PairedCharacterConditionAlt}(C,C_T)$ on a good cut.
In particular the two type primes are distinct, $p=g_k$, the
size, sign and host-jump hypotheses hold, and the specified
rank Selmer and cyclotomic lengths are finite. Write $n=r+t$,
$f=[E^+:\mathbb Q]$. Denote by $A_\pi$ the part of the
native right hand side independent of the table constant and
of the quadratic reserve:
\[
 A_\pi=s_\pi+\sum_{0<j<t}\mathrm{shaLen}_p(j)
        +v_p(h_E)+4n v_\delta+t\sum_{q\in Q_i}v_\zeta(q).
\]
Suppose that at each $\pi\in Y$ the length from P84D1 satisfies
the native point bound, namely there are integers $2\le N\le n+1$,
$1\le u$, $4u\le N^2$, for which
\[
 c_\pi\le 4 f n C_T\bigl(Nu+(\eta+t+1)(N+u)\bigr)+A_\pi .
\]
The witnesses may vary with $\pi$. For every $C_T\in\mathbb N$
this implies the paired inequality with $C=C_T$.
Indeed the native rank assumptions imply $n\ge3$. Then
$N\le4n/3$, $u\le4n^2/9$, $N+u\le8n^2/9$ and
$Nu\le16n^3/27$. Since $N+u\le 8n^3/27$ as well, we have
\[
 Nu+(\eta+t+1)(N+u)
 \le \tfrac{24}{27}n^3+\tfrac89 n^2(\eta+t)
 \le n^2(n+\eta+t).
\]
The native right hand side adds $2fr^2$ to
$4fnC_Tn^2(n+\eta+t)+A_\pi$. Multiply by the
nonnegative $m_\pi$, sum, and use P84D1. This gives the
precise conclusion for the tested families and parameters.
If the point bound holds for every native test, the paired
condition follows with the same uniform constant. Neither
the argument nor the alternate test requires an extra
$\mathrm{SetSQ}$ assumption. P84D2 also proves the same
conclusion directly whenever its Smith differences divided by
$e$ are bounded by the native right hand side at every test.
% END P84D3

% BEGIN P84D4
\paragraph{P84D4. A cut carrying one prescribed tuple of local types.}
Work on the good cut of P84B1, writing $X$ for its spectrum
and $Y=Z=Z_{\rm good}$. Fix a type-prime index and $\pi\in Y$.
Its good list supplies a parameter $\sigma$ whose polynomial
labels form one tuple $a$. On every representative above $\pi$,
the same-list comparison of P84C2 puts its whole type in
$\mathrm{MtPar}(\sigma)$. By the proof of P84C2, $E(a)$ fixes
every such vector. Use just the one integral central element
lifting $N_G E(a)$ from P84C1 and let $\beta$ be its action
on the cut lattice. It commutes with Hecke. The transported
$\mathrm{EigSpan}$ comparison gives $\beta=N_G$ on the entire
rational eigenspace at $\pi$.

Take $S_\pi$ to consist of those $x$ admitting a representative
and a type vector on which $E(a)$ acts nontrivially. This
contains $\pi$, since a representative with nonzero type exists.
At every $x\notin S_\pi$ all representatives have zero action,
so $\beta$ vanishes on the whole rational eigenspace there.
Also $S_\pi\cap Z\subseteq Z_{\rm good}$ since this is a good
cut. The integer $b=N_G$ is nonzero and
$\ell_O(O/bO)/e=v_p(N_G)$ by factoring it in the coefficient
DVR. These are precisely the conditions of
$\mathrm{CutAt}(i,k,\pi,S_\pi,\beta,b)$.

For $x\in S_\pi$ retain the same witnessing representative
and vector. Each of the commuting local series factors of
$E(a)$ acts nontrivially on that vector. Expanding a local
series into its finite sum of irreducible character projectors
shows that some one of these projectors acts nontrivially.
The local type occurrence identity in the given depth-zero
record then puts that irreducible in the representative's
specified set of types. Its label is the prescribed polynomial
at that place. At places over $g_k$ these polynomials are
squarefree and avoid both $\chi_J$ and its inverse by P84C3.
Thus the single tuple of polynomials, followed at each $x$ by
one representative before all places, witnesses the full
$\mathrm{TypeLocus}(P,i,k,S_\pi)$. The set $S_\pi$ is chosen
before the good Galois family or Selmer data tested at $\pi$.
% END P84D4

% BEGIN P84D5
\paragraph{P84D5. The cost of choosing a different local type at each point.}
For each $\pi\in Y$ use P84D4's cut or any native
$\mathrm{CutAt}(i,k,\pi,S_\pi,\beta_\pi,b_\pi)$. Put
$B_\pi=S_\pi\setminus Z$. Set $c_\pi^*=
[d(B_\pi\cup\{\pi\})-d(B_\pi)]/e$ from the algebra on
this spectrum. P84D2 identifies it with
$\mathrm{vccCut}(i,S_\pi,\pi)$, since only vanishing on
$B_\pi$ is tested. Then
\[
 \eta_{\rm loc}\le\sum_{\pi\in Y}m_\pi
            \bigl(c_\pi^*+v_p(N_G)\bigr).
\]
To see this even when all the supports are different, use LN4's
filtration of the projected lattice $U=P_Y M$ modulo
$M\cap V_Y$, where $V_Y$ is the sum of the rational eigenspaces
in $Y$. At its step for $\pi$, choose $h\in T$ vanishing on
$B_\pi$ whose value at $\pi$ generates the evaluation ideal.
Existence follows from P84D1 with this vanishing set. The
endomorphism $h\beta_\pi$ is integral, kills every rational
eigenspace outside $Y$, and acts at $\pi$ by $h(\pi)b_\pi$.
Indeed $\beta_\pi$ preserves each eigenspace by commutation,
and outside $Y$ either it vanishes or $h$ does. Therefore
$h\beta_\pi$ sends $U$ into $M\cap V_Y$: for $P_Ym\in U$ its
image equals $h\beta_\pi m$. It preserves the partial sums of
eigenspaces in the filtration. On the current graded piece
it acts by $h(\pi)b_\pi$, which consequently kills the
graded quotient. That quotient is generated by at most $m_\pi$
elements by the free-lattice argument of LN4. Summing its
length bounds and using the scalar length test in
$\mathrm{CutAt}$ proves the inequality.

At a paired test let $g=p$ and $g'$ be the other generator
prime. The logarithm guard gives $g'\le g^2$. There are $f$
factors at each prime in $N_G$. The elementary column count for
$\mathrm{GL}_r(\mathbb F_q)$ gives its order as
$q^{r(r-1)/2}\prod_{a=1}^r(q^a-1)$. For the $g$ factors
only the initial power contributes. For $g'\ne g$, each
positive integer $g'^a-1<g^{2a}$ has $g$-valuation at most
$2a$. Hence
\[
 v_g(N_G)\le f\{r(r-1)/2+r(r+1)\}\le 2fr^2.
\]
This is the quadratic reserve in P84D3. Apply its polynomial
estimate to $c_\pi^*$ and absorb this scalar cost before summing
multiplicities. The native point bounds on these smaller
$S_\pi$ now suffice for $\mathrm{PairedCharacterConditionAlt}$
with $C=C_T$. In particular P84B2 applies with these
pointwise tests in place of the tests on the whole retained
spectrum. The Smith differences for the larger bad set need
not satisfy the smaller-set bounds.
% END P84D5

\section{Finite cocycles and the two type-prime local bounds}

% BEGIN T84L1
\paragraph{T84L1. The lifting exponent on a family over a degree-one place.}
Put $p=\ell_i$, $N=n_V+1$, $a=\eta+t+1$. The degree-one test on
the fixed completion identifies $K=E_w$ with $\mathbb Q_p$ inside
the given algebraic closure. Suppose $a<H$ and
$C_T(HN^2+N^3)<p$. In $\mathrm{PotFam}$ take its common finite
normal extension $K'/K$, its Frobenius $F$ and its $F$-stable
set $\Sigma$ of at most $N$ inertia characters containing $1$.
The crystalline period action on the unramified constants of $K'$
is the one in that datum. In $\mathrm{JumpFam}$ take the common
set $J\subset[-a,H]$, containing $0$ and of size at most $N+1$.

The lattice maps tested by $\mathrm{LiftFamW}$ can have arbitrarily
large source and target. Their rational spaces embed equivariantly
in finite sums of family members. After restriction to $K'$ each
is crystalline by the subobject and direct-sum properties of the
Fontaine comparison. Their jump sets at $K$ stay in $J$: tensor
an injection into the sum with $B_{\rm dR}$, and choose an
ordinary linear left inverse of the injection. The induced map
on the period tensor spaces then pulls the degree-$d$ filtration
back to the degree-$d$ filtration for every $d$. Taking Galois
invariants preserves this intersection. A nonzero graded
piece thus injects into the corresponding graded piece of the sum.
Similarly on $D_{\rm cris,K'}$ the injection respects inertia.
Each character in $\Sigma$ is trivial on the indicated subgroup
over $K'$ and has order prime to $p$, by $\mathrm{InertiaThrough}$
at any one family member. On the ambient period space the
action factors through a finite abelian quotient, the image of
these characters. Its character projectors preserve the image
of the injected period space, which is inertia stable by
normality. Thus the same $\Sigma$ works for each tested
subrepresentation. For an empty family all tested spaces are
zero and the lifting assertion is immediate.

Let $r'$ be the largest Frobenius orbit length of
$\alpha\beta^{-1}$ for $\alpha,\beta\in\Sigma$. Frobenius
permutes $\Sigma$ and fixes the trivial character. For two
characters in the same orbit, its length bounds that of their
ratio, and is at most $N-1$ unless the orbit is the trivial
singleton. For two distinct orbits of lengths $d_1,d_2$, the
ratio is fixed after $\operatorname{lcm}(d_1,d_2)$ steps, and
$d_1d_2\le(d_1+d_2)^2/4\le N^2/4$. Since $N\ge2$, this
gives $r'\le\lfloor N^2/4\rfloor$ in all cases. Now apply
the stated $\mathrm{admT}(C_T)$, the quantitative potentially
crystalline lifting theorem, to the two tested lattices and
their surjection. Its hypotheses on these very spaces have
just been verified, including the full labelled jump set.
It gives lifting with exponent $C_T(Nr'+a(N+r'))$. Multiply
the finite lift and its integral coboundary by
$p^{e_{\max}-C_T(Nr'+a(N+r'))}$, where
\[
 e_{\max}=C_T\bigl(N\lfloor N^2/4\rfloor+
                         a(N+\lfloor N^2/4\rfloor)\bigr).
\]
This proves $\mathrm{StepLift}$ with its original constant.
The quantitative lifting theorem is an explicit input here.
% END T84L1

% BEGIN T84L2
\paragraph{T84L2. A cut cocycle as the pushout of an integral rational coboundary.}
Write $L$ for the local coefficient field and $O$ for its
integers. Let $M$ have rank $k$ and satisfy $\mathrm{SubTw}$.
Thus $M$ is continuous and its rational realization embeds
globally in one of the retained twists $r_x\chi_1^{-1}$.
Write $U$ for its local rational space and $T$ for its standard
lattice. A generator of $\mathrm{cutCoc}$ is given, over the
original coefficients, by a finite tuple of retained twists,
a finitely generated stable module $Y$ in their sum, a vector
$E$ with $y(g)=gE-E\in Y$, and an equivariant integral linear
map $f:Y\to U/T$ under the coefficient identifications.

There is a global stable complement to the rational span $X$ of
$Y$. Indeed each good rational representation is semisimple,
twisting by a nonzero character preserves every invariant
subspace and its complements, and a finite direct sum of
semisimple modules is semisimple. Apply the equivariant projection
onto $X$ to $E$. Then $y(g)=gE_X-E_X$ since the projection
fixes every $y(g)$. Transport this decomposition to the
completion and restrict to $G_K$. The actual image of $Y$ under
the given coefficient isomorphism is a stable full lattice in
$X$, being finitely generated, torsion-free and spanning. It
retains the map $f$. Thus $y$ is a continuous integral cocycle
whose rational class in $X$ is zero. In particular it is in the
integral Bloch--Kato finite condition. Its pushout $f(y)$ is
the original local function on $G_K$ even when $E_X\notin Y$.

The same global complement construction applies to $U$ inside
its chosen member. It follows that both $X$ and $U$ embed in
finite local direct sums of the twisted family, and that their
Tate duals $X'=X^*(1)$ and $U'=U^*(1)$ embed in corresponding
sums of the dual family: extend functionals by zero on the
transported complements. The pairings transform by the
cyclotomic character. In either direction all the objects in
the graph construction below therefore pass the rational
embedding tests of $\mathrm{LiftFamW}$. This argument uses
semisimplicity before restriction to a decomposition group.
The common $K'$ of the direct potential family makes $X,U$
and their Tate duals crystalline over $K'$ by Fontaine's
subobject, sum and dual properties. Potentially crystalline
implies de Rham over $K$.
% END T84L2

% BEGIN T84L3
\paragraph{T84L3. The direct and transposed lattices.}
For the direct lifting condition at exponent $e$, choose an
$L$-linear lift $s:X\to U$ of $f$ on $Y$, by lifting the
images of an $O$-basis. In $X\oplus U$ take the lattice
\[
          B=\{(v,u):v\in Y,\ u-sv\in T\}.
\]
A basis is obtained from the graph of the chosen basis of $Y$
together with a basis of $T$. Equivariance of $f$ says exactly
that $gsv-sgv\in T$ for $v\in Y$, so $B$ is stable.
Projection onto $Y$ is surjective. Apply the finite-cocycle
lifting condition of the direct family to $y$. We obtain a
finite $B$-valued cocycle $(b_1,b_2)$ and $v\in Y$ with
$p^ey=b_1+\delta v$, where $\delta v(g)=gv-v$.
Thus $p^ef(y)=\overline{b_2}+\delta f(v)$. Projection takes a
rationally finite cocycle to a finite one, so this lies in
$H_{f,D}(T)$, the reduction of rational finite cocycles plus
all coboundaries with values in $U/T$.

If only the dual family lifts at $e$, put
$Y'=Y^*(1)$, $T'=T^*(1)$ and let $s^*:U'\to X'$ be the
transpose under the perfect pairings with values in $L(1)$.
In $U'\oplus X'$ use
\[
          B'=\{(b,c):b\in T',\ c-s^*b\in Y'\}.
\]
This projects onto $T'$ and is stable. Pair the transpose's
equivariance defect with $v\in Y$ and use the preceding
defect condition for $g^{-1}v$; the cyclotomic factor
preserves $O(1)$. Given any integral finite test cocycle
$d$ with values in $T'$, dual lifting gives a finite
$B'$-valued cocycle $(b,c)$ with $p^ed=b+\delta t'$ for
some $t'\in T'$. In the Tate cup of $b$ with $f(y)$, the
representative at $(g,h)$ is
$\langle g s(y(h)),b(g)\rangle$ modulo $O(1)$.
Its difference from $\langle g y(h),c(g)\rangle$ is integral:
the first error pairs $g s(y(h))-s(g y(h))\in T$ with $b(g)\in T'$,
and the second pairs $g y(h)\in Y$ with $c(g)-s^*b(g)\in Y'$.
The latter rational cup is a coboundary since $y=\delta E_X$
rationally. Reducing its continuous rational primitive modulo
$O(1)$ preserves that identity. Hence the Tate pairing of
$p^ef(y)$ with $d$ vanishes, since an integral coboundary in
the $T'$ argument has zero pairing. Integral Bloch--Kato
orthogonality for the de Rham representation $U$ identifies
the annihilator of all these tests with $H_{f,D}(T)$.
This proves $p^ef(y)\in H_{f,D}(T)$ on the dual side as well.
% END T84L3

% BEGIN T84L4
\paragraph{T84L4. The rank term and the full cut span.}
Let $A=H^1_f(K,T')$, the inverse image of the rational finite
condition in integral cohomology. Local finiteness makes $A$
finitely generated over $O$. Its rank is $\dim_L H^1_f(K,U')$:
a continuous cocycle on the compact group $G_K$ has bounded
image, so some nonzero scalar makes it integral. The same
denominator clearing applies to coboundaries. Over
$K=\mathbb Q_p$ the Bloch--Kato dimension formula gives
\[
 \dim_L H^1_f(K,U')
 =\dim_L (U')^{G_K}+\dim_L D_{\rm dR}(U')
                         -\dim_L {\rm Fil}^0D_{\rm dR}(U')\le k .
\]
Indeed $v\mapsto v\otimes1$ embeds the invariants in the
zero filtration, and $\dim_L D_{\rm dR}(U')=k$ by de Rhamness
and the degree-one hypothesis.

The torsion subgroup of $A$ needs at most $k$ generators too.
It lies in the kernel of $H^1(K,T')\to H^1(K,U')$, the image
of $(U'/T')^{G_K}$ by the connecting map. Take preimages of
finitely many generators of the torsion in this image. Their
span is a finite submodule of $(L/O)^k$, since their
denominators are jointly bounded. Such a submodule needs at
most $k$ generators by elementary divisors over a DVR. The
images of its generators generate the given torsion subgroup.
Consequently
\[
   \operatorname{length}_O(A/p^e A)\le
                2ke\operatorname{length}_O(O/pO).
\]
This also covers $e=0$ and $k=0$.

For the entire $O$-span $S$ of cut cocycles, Tate evaluation
now sends $S/(S\cap H_{f,D}(T))$ injectively into
$\operatorname{Hom}_O(A/p^e A,L/O)$. The zero kernel is precisely
Bloch--Kato orthogonality. The factorization through $A/p^e A$
holds for every generator by T84L3 and hence for the span.
The $O$-linear invariant map for $O(1)$ gives the scalar target
$L/O$ (equivalently choose a basis for the inverse different
in the scalar trace pairing). A finite $O$-module and its
$L/O$-dual have equal length, as one checks for $O/(\varpi^b)$
and then sums. The quotient on the left is
$(S+H_{f,D}(T))/H_{f,D}(T)$. This proves
$\mathrm{LocBoundP}(2ke)$ simultaneously for all test
lattices, without adding one cost per tuple or per point.
In particular both $\mathrm{StepLocDirect}$ and
$\mathrm{StepLocDual}$ of the type-prime cut follow from their
respective finite-cocycle lift and the assumed local
Fontaine, Bloch--Kato and Tate comparisons. When T84L1
provides that lift, $e=e_{\max}$. The production of the common
potential family and of the appropriate signed jump family,
and the quantitative hypothesis $\mathrm{admT}(C_T)$ used in
T84L1, are separate inputs.
% END T84L4

% BEGIN T84L5
\paragraph{T84L5. The two signs of the host jump condition.}
At $w\mid p$ put $c_0=(t-1)/2$, $H=s+S$, $a=\eta+t+1$
with $a<s$. The host test bounds $|U_w|\le n_V$ and contains
the jumps of every direct member at every embedding. Its
integer entries $b$ either lie between $0$ and $t-1$, satisfy
$|b-c_0|\le\eta+1$, or lie in the specified one-sided interval.
For the positive sign this last interval is
$c_0<b\le H$. All three alternatives lie in $[-a,H]$ since
$0\le c_0\le t-1$ and $a<s$. Thus $J=U_w\cup\{0\}$
suffices for every retained point.

For the negative sign the one-sided interval is
$2c_0-H\le b<c_0$. Reflect by $b\mapsto-b-1$. On this
interval the reflected value is at most $H-2c_0-1\le H$
and at least $-c_0-1\ge-a$. For $0\le b<t$ the reflected
value is in $[-t,-1]$. On the central interval it lies in
$[-c_0-\eta-2,\eta-c_0]$, again contained in $[-a,H]$
since $c_0+1\le t$. Take $J=\{-b-1:b\in U_w\}\cup\{0\}$.
Both chosen sets have size at most $n_V+1$, within the native
$N+1$ allowance. For an actual direct potential family each
point is de Rham by descent from $K'$. The dual matrix at $g$
is $\epsilon(g)R(g^{-1})^{\mathsf t}$ for $R=r_x\chi_1^{-1}$.
The dot-product pairing is perfect with values in $L(1)$;
the dual filtration reverses each labelled grading and $L(1)$
translates it by $-1$. Therefore a dual jump at $b'$ corresponds
exactly to a direct jump at $-b'-1$. This proves the two
$\mathrm{JumpFam}$ conclusions from the one host test.

The very same $K',F$ supply the dual potential family. Over $K'$
the Tate dual is crystalline. On crystalline periods the
cyclotomic line has basis $e_1\otimes t_{\rm cris}^{-1}$,
fixed under descent since the two cyclotomic factors cancel.
By the crystalline tensor and dual comparison, the pairing
on periods is perfect over $K'_0\otimes_{\mathbb Q_p}L$ and its
inertia action on that line is trivial (inertia fixes $K'_0$).
A component dual to the $\alpha$-space therefore transforms by
$\alpha^{-1}$. Inverting all characters preserves the
trivial character, their number, their restrictions over
$K'$, and their prime-to-$p$ orders; inversion commutes with
Frobenius conjugation. Thus $\Sigma^{-1}$ supplies
$\mathrm{PotFam}$ on the dual side without a new field or a
new size bound. Direct and dual matrix functions are continuous
by continuity of $r_x$, the characters, inversion in $G_K$ and
transpose. This supplies the sign-specific hypotheses of T84L1
from the direct potential family and the stated host test under
the native period comparisons.
% END T84L5

% BEGIN T84L6
\paragraph{T84L6. The rank block and each cyclotomic line at a type place.}
Let $\pi$ belong to the retained set, and use the given rational
semisimple decomposition of $r_\pi$ into the rank-$r$ block
$\rho_\pi$ and the characters $\chi_1,\ldots,\chi_t$. Twist
the same isomorphism by $\chi_1^{-1}$. Its inverse on the first
factor embeds $W_{\pi,O}$ rationally; on the $(c-1)$st
zero-based character coordinate it embeds
$\chi_c\chi_1^{-1}$ for every $1\le c\le t$. Each entry
of $\rho_\pi(g)$ is a fixed rational linear combination of the
entries of $r_\pi(g)$. Clearing its denominator gives an
integral continuous function; multiplication by a nonzero
element of a DVR is a homeomorphism onto its image in the
adic topology. The rank block therefore is continuous, as are
its twist and the character ratios. These embeddings give the
full $\mathrm{StepSub}$ tests for T84L2.

Set $q=c-1$. The original character formula on this range gives
$\chi_c\chi_1^{-1}=\epsilon^{-q}$, and the local coefficient
bridge identifies $\epsilon$ with the period cyclotomic
character. Under a coordinate basis the local lattice is thus
$O(-q)\subset L(-q)$. For the odd type prime $p$ and its
degree-one completion, if $q>0$ then the rational invariants
vanish (the cyclotomic image is infinite), and the de Rham
filtration has its unique step at $q>0$, so its degree-zero
quotient $D_{\rm dR}/{\rm Fil}^0$ is zero. The dimension formula
then gives $H^1_f(K,L(-q))=0$. Thus every finite rational
cocycle is a coboundary, and after reduction modulo $O(-q)$
the native finite cocycles are precisely the coboundaries.

For $q=0$ the same formula gives dimension one. The unramified
degree homomorphism is nonzero and crystalline: in the unramified
Witt periods solve $\varphi b-b=1$ by first solving
$z^p-z=1$ over the residue algebraic closure, then correcting
the error successively modulo higher powers of $p$ by the
same surjectivity of $z\mapsto z^p-z$. Inertia fixes this
period and $\varphi^m b-b=m$ on integral powers of arithmetic
Frobenius, hence on their closure by continuity. This spans
the rational finite condition. Reduction modulo $O$ therefore
vanishes on inertia. Conversely an unramified $L/O$-valued
homomorphism lifts on Frobenius to some $d\in L$, and the
continuous map $\widehat{\mathbb Z}\to\mathbb Z_p\to L$,
the last map being multiplication by $d$, gives the lift.
The native inertia identification gives the claimed unramified
test on the specified local bridge. This proves
$\mathrm{StepHfChar}$ and, with T84L1--T84L5, supplies the
rank-$r$ and all rank-one local estimates with exponents
$2r e_{\max}$ and $2 e_{\max}$ at every tested type place
once the direct potential family and $\mathrm{admT}(C_T)$
are given with the host, size and comparison hypotheses.
% END T84L6

\section{A common potentially crystalline family at a type prime}

% BEGIN K84P1
\paragraph{K84P1. The local comparison on the two places of a type label.}
Use the native hypotheses of $\mathrm{StepTypeAlt}$, including
$\mathrm{admDZ}(P)$, the three automorphic and Galois inputs
$\mathrm{adm1},\mathrm{adm2}_{\rm alt},\mathrm{adm5}$ and the
comparison laws $\mathrm{tie0}$. Fix $\ell_i=g_k=p$, an automorphic
representative above $x$, and an occurring type with squarefree label
$s(X)$ avoiding both designated units. The depth-zero test specifies
one of the two units $u$ and gives at the distinguished split place
the degree $n_V=r+t$ label $s(X)(X-u)^t$, with zero monodromy.
At its conjugate place it gives the multiset of inverses of these
roots, again with zero monodromy and full degree. The constant term
is nonzero by the occurring-type label condition. The calculation of
P84C3 applies to this individual label: the tame generator
distinguishes the nonzero root characters, and only $u$ (respectively
$u^{-1}$) can occur twice. Its order divides $p-1$. This gives the
noncritical condition $\mathrm{Nc}$ at both places. No goodness
condition on other points is needed.

Choose a base-change list for this representative by $\mathrm{adm1}$.
At the tested split place its local sum is the indicated parameter.
Noncriticality passes to each summand. Wild and zero-monodromy
conditions pass by direct sum, inertia polynomials multiply, and
the repetition threshold in the definition is nonincreasing with
dimension. An eigenvalue reaching its summand's threshold therefore
reaches the total threshold. The real unramified twists preserve
all these conditions. The centered cuspidal realizations of
$\mathrm{adm2}_{\rm alt}$, with the algebraic twists required
by regularity of the list as in C84F5, now give a de Rham sum. More explicitly, the top
member of each length-$d$ string has its infinity weights in the
regular total weight shifted by $(d-1)/2$. This fixes the parity
of the algebraic twist. Undo the common centering and take the
$d$ successive Tate shifts of this one realization; direct sums
and these twists preserve de Rhamness.
On split unramified Frobenius elements its polynomial agrees with
that of the good $r_x$, because both are the normalized Hecke
polynomial of the given point. Density and semisimplicity identify
the two representations over the algebraic closure of the
coefficient field, so the de Rham and Weil--Deligne comparisons
transfer by $\mathrm{tie0}$.

For clarity, inertia agreement in the weakened cuspidal comparison
has a guard. On a commutator of inertia, each nonzero-root character
of the label is one. The full product polynomial there is therefore
$(X-1)^{n_V}$, and the corresponding polynomial for each summand
is $(X-1)$ to that summand's dimension, by monic divisibility.
Unramified normalizing twists have the same inertia polynomials.
Together with zero monodromy these facts meet the commutator
guard in the comparison. In degree one the comparison is the
character identity. On the other summands, apply the guarded
inertia identity and multiply over the direct sum and its Tate
shifts. Thus at either place $w$ the Galois parameter has zero
monodromy and for every $h\in I_w$ its inertia polynomial is
\[
              \prod_{y\in M}(X-\omega_{k,w}(y,h)),
\]
where $M$ is the specified root multiset, inverted at the conjugate
place. The same chosen $u$ serves both places. This is
$\mathrm{TypeLGC}(P)$ with the required quantifiers.
% END K84P1

% BEGIN K84P2
\paragraph{K84P2. The common field and the bounded set of tame labels.}
Assume now the native premises of $\mathrm{StepPotDirect}$, in
particular $\mathrm{TypeLGC}(P)$, the local bridge and character
ties, the crystalline restriction and tensor comparisons and
$\mathrm{TieInePoly}$. Fix its $i,k,\pi,X_s,w$.
The type locus has a single tuple of squarefree labels. Since
$\pi\in X_s$, each local label occurs and hence is monic of
degree $r$, with nonzero constant term. At the distinguished
member of the pair let $c$ be the designated residue unit, and use
\[
           Y=\{\text{roots of }s\}\cup\{c,c^{-1}\}.
\]
At the conjugate place invert this set. The choice of $u$ in
K84P1 may vary with the point, but every resulting multiset is
supported in $Y$. The separate character tie supplies the label
$u_0\in Y$ of the actual de Rham line $\chi_1$ at this $w$,
including inversion at the conjugate place. Its nilpotent
monodromy is zero by dimension one. Thus each $r_x$ and the line
are potentially crystalline, by the de Rham and zero-monodromy
comparison transported through the coefficient bridge.
The full inertia polynomial on the line is
$X-\omega_{k,w}(u_0,h)$ by the one-dimensional label identity.

Write $m_G$ for the product group's order in the datum.
Every $y\in Y$ satisfies $y^{m_G}=1$: a root of degree
$d\le r$ has order dividing $p^d-1$, and
$|\mathrm{GL}_r(\mathbb F_p)|=\prod_{j=0}^{r-1}(p^r-p^j)$
contains the factor $p^{r-d}(p^d-1)$. The designated units have
$d=1$, and inversion preserves the assertion. Hence all their
orders are prime to $p$. Moreover $|Y|\le r+2\le n_V+1$.

The type prime splits completely, so the completion $K$ has
degree one over $\mathbb Q_p$. For each point take a finite
extension over which its representation becomes crystalline; do
likewise for $\chi_1$. For each $y\in Y$ the open-kernel clause of
the tame label supplies an open subgroup on whose inertia its
character is trivial. Via the decomposition-group bridge this
contains the inertia of a finite local extension: an open
neighbourhood of the identity in inertia contains the intersection
with an open subgroup of the whole local Galois group. Take the
normal closure $K'$ of the compositum of these fields and $K$.
There are finitely many points and labels. Crystallinity persists
over $K'$ by the given restriction comparison. Normality makes
$K'_0=K'\cap\mathbb Q_p^{\rm nr}$ stable, with the prescribed
action on unramified constants. The same $K'$ works on all points
and on the one-dimensional character, and it kills the inertia
characters attached to $Y$ on $I_K\cap G_{K'}$.
% END K84P2

% BEGIN K84P3
\paragraph{K84P3. Recovering simultaneous characters on crystalline periods.}
Put $H=G_{K'}$ and $I=\ker(G_K\to\mathrm{Aut}(K'_0))$.
Then $I/H$ is the image of $I_K$. Indeed the fixed field in $K'$ of the image of $I_K$ is
$K'\cap K^{\rm nr}=K'_0$, since $K=\mathbb Q_p$.
The comparison on period-ring constants identifies this $I$ with
the inertia group of the specified period datum. The inertia
identification of the bridge reads $I_K$ as the designated
global inertia.

Lift $\omega_{k,w}(y,h)$ to the coefficient DVR by the supplied
primitive $m_G$th root and then map it into the local coefficient
field $L$. This lift exists and is unique: the value is an
$m_G$th root by multiplicativity in $y$, and the powers of the
chosen primitive root have distinct complex images exhausting
those roots. The two embeddings of the DVR are the given ones.
The lifts are multiplicative in $y$ and in $h$. On the cyclic
group generated by $Y$ they kill $I_K\cap H$ by K84P2. They
therefore descend uniquely to characters $a_y:I/H\to L^\times$.
Each has order dividing the order of $y$. Moreover all are
powers of one character: take a generator of the finite cyclic
group generated by $Y$ in $\overline{\mathbb F}_p^\times$,
and apply the same multiplicative lift.

For any one of our representations let $M$ be its specified
multiset and $D'=D_{\rm cris,K'}$. The action of $I$ on $D'$
is $L$-linear and factors through the finite group $I/H$.
For each $g\in I$ choose $h\in I_K$ in its coset modulo $H$.
The polynomial with roots the lifted $\omega_{k,w}(y,h)$,
$y\in M$, has coefficients in the original DVR and its complex
image is the prescribed inertia polynomial. By the
$\mathrm{TieInePoly}$ premise the local image of this polynomial
annihilates every crystalline period under $h$. On these periods
$h$ and $g$ act identically. Thus the action of $g$ is annihilated
by $\prod_{y\in M}(X-a_y(g))$, with the multiplicities in $M$.

Let $C$ be the cyclic image of the single character generating
the $a_y$. An element in its kernel has unipotent action by the
last polynomial. Since it has finite order on the period space
in characteristic zero, its action is the identity. On $D'$
the action now factors through $C$. Choose a preimage of a
generator. Its minimal polynomial divides both the above
product and the separable polynomial $X^{|C|}-1$; hence it is
diagonalisable over $L$ with eigenvalues among the listed
$a_y$ evaluated at that generator. Every eigenspace is a
simultaneous character space for the entire cyclic quotient.
This proves the full period eigenspace sum, the triviality on
$H$, and the prime-to-$p$ orders required by
$\mathrm{InertiaThrough}$. For $\chi_1$ the allowed set is
the singleton $\{a_{u_0}\}$ throughout.
% END K84P3

% BEGIN K84P4
\paragraph{K84P4. A Frobenius-stable support for the twisted family.}
By the crystalline tensor and dual comparison over the common
$K'$, $r_x\chi_1^{-1}$ is crystalline there, and every character
of its period inertia action belongs to
$S=\{a_ya_{u_0}^{-1}:y\in Y\}$. The evaluation pairing of the
line with its inverse divides the inertia eigenvalue by
$a_{u_0}(g)$ for every $g$. Thus $1\in S$ and $|S|\le n_V+1$.

Take a lift $F\in G_K$ of the required residue-field Frobenius on
$K'_0$. It exists because $K=\mathbb Q_p$ and the automorphism of
this unramified field extends to the algebraic closure. The action
of $F$ on $D_{\rm cris,K'}$ preserves periods by normality
and is $L$-linear, with its specified semilinear action on $K'_0$.
For a nonzero vector in the $a$-space, its image under $F$ belongs
to the character $g\mapsto a(F^{-1}gF)$, since $I$ is normal.
It is nonzero and therefore this character occurs in the same
representation and belongs to $S$. Take $\Sigma$ to be $1$
together with the characters that actually occur at any retained
point. The preceding argument makes $\Sigma$ Frobenius-stable,
and K84P3 shows that it gives every required inertia decomposition.
As $\Sigma\subset S$, stability has incurred no cardinality loss.
These $K',F,\Sigma$ give $\mathrm{PotFam}$ for the whole direct
family. This proves $\mathrm{StepPotDirect}$ from its stated
inputs. T84L5 supplies the Tate-dual family with the same field,
and provides the two signed jump tests under the host hypothesis.
The quantitative lifting theorem $\mathrm{admT}(C_T)$ enters
later in T84L1.
% END K84P4

\section{A common tame field and the constant in the type-prime lifting theorem}

% BEGIN Q84T1
\paragraph{Q84T1. Integral descent in the coefficient field.}
Take any of the local data $D$ tied to the usual periods and any test
surjection $g:T_1\twoheadrightarrow T_2$ of stable free lattices for
$G_{\mathbb Q_p}$ in the definition of $\mathrm{admT}$. Its potentially
crystalline representations use a common finite normal extension $K'$,
an inertia type $\Sigma$ and Frobenius as in T84L1. Write
$N\ge2$, $1\le b<H$, $0\in U\subset[-b,H]$, $|U|\le N+1$ for
the jump parameters. Assume $40(HN^2+N^3)<p$. Put $a=(p-1)!$,
$e=p^a-1$, and enlarge the coefficient field $L$ to its compositum
$L'$ with the unramified extension $\mathbb Q_{p^a}$ of degree $a$.
Its residue field contains $\mathbb F_{p^a}$, so its Teichm\"uller units
contain $\mu_e$. Both $a$ and $e$ are prime to $p$.

Write $O,O'$ for the integer rings. If $0\ne d\in O$, $x\in O'$
and $dx\in O$, then $x\in L\cap O'=O$. Thus $O'/O$ is finite
torsion free over the DVR $O$, hence free. There is an $O$-linear
retraction $s:O'\to O$ with $s(1)=1$. Scalar extension preserves
the surjection, its kernel, its periodic characters and its jumps.
For the latter assertion, choose an $L$-basis of $L'$: the period
invariants and the intersections with each filtered period space
are computed coordinate by coordinate. For a finite integral
cocycle $c$ on $T_2$, an identity
$p^d(1\otimes c)=g_*c'+\partial t'$ after extension descends
by $s\otimes1$. The two images are integral, and the extension
of $s$ to $L'\to L$ takes a crystalline period primitive for
$c'$ to one for its image. This proves descent of the lifting
exponent without a factor $[L':L]$.
% END Q84T1

% BEGIN Q84T2
\paragraph{Q84T2. A common tame field for every input lattice.}
Choose $\pi_0^e=-p$ and $\pi_{j+1}^p=\pi_j$, and put
$K_t=\mathbb Q_{p^a}(\pi_0)$. The ordinary inertia group $I$ maps
onto the inertia group of the finite normal extension $K'$.
Indeed a lift of an element acting trivially on the residue field
of $K'$ can be corrected by an element of $G_{K'}$ with the inverse
residue action. Thus characters of the inertia group in the period
datum that are trivial on $G_{K'}$ are distinguished by their
restrictions to $I$. Their restrictions to wild inertia are trivial:
the finite quotient is a $p$-group there and each character has
order prime to $p$. On tame inertia absolute arithmetic Frobenius
acts by $p$-th power. The selected period Frobenius agrees with it
on the unramified constants of $K'$, so their difference lies in
$I G_{K'}$. Conjugation by this difference fixes every character
of that group trivial on $G_{K'}$. The period of each character is at most
$|\Sigma|\le N<p$, hence divides $a$. Its values therefore have order
dividing $e$. The map $\omega(g)=g(\pi_0)/\pi_0$ identifies the
cyclic quotient of tame inertia of order $e$ with $\mu_e$:
$X^e+p$ is Eisenstein over $\mathbb Q_{p^a}$ and splits after adjoining
$\pi_0$. Consequently the restricted characters are uniquely
$\omega^c$ for $c$ in a set $A\subset\mathbb Z/e$. We have $0\in A$,
$pA=A$, and $|A|\le N$. For $R=R(A)$, the largest period under $p$
of $c-d$ with $c,d\in A$, injectivity on restricted characters
identifies $R$ with the original ratio period. For two different
orbits of sizes $i,j$, the period divides $\mathrm{lcm}(i,j)\le
ij\le N^2/4$. Within the same orbit it divides its length, at most
$N-1$ outside the fixed point $0$. Thus $1\le R$ and $4R\le N^2$.

For either representation $V$, pass to $K''=K_tK'$. The crystalline
period module there is obtained from that over $K'$ by extending
unramified constants (the old periods already supply full dimension).
Inertia of $K''/K_t$ fixes those old periods by the character
decomposition and $\omega|_{I\cap G_{K_t}}=1$, and fixes the new
unramified constants. The remaining action of
$\mathrm{Gal}(K''/K_t)$ is semilinear through its residue-field
quotient. Galois descent for vector spaces gives a full-dimensional
space of crystalline periods over $K_t$. For $g\in I$ the operator
$\prod_{c\in A}(g-\omega(g)^c)$ kills the old periods by their
decomposition, hence the new periods by scalar extension and restriction.
The kernel lattice $T_0=\ker g$ is free and stable. Its rational
space inherits crystallinity, the annihilating polynomial and jumps in
$U$ by crystalline subobject exactness and strictness of de Rham
filtrations, as in T84L1. This produces the three good lattices for
the same $A,U,K_t$.
% END Q84T2

% BEGIN Q84T3
\paragraph{Q84T3. Exactness at the middle of a realized integral triple.}
Suppose that compatible Breuil--Kisin realizations of the three
lattices and the two maps are given on $K_t$. Write $F,G$ for their
matrices over $\prod_{\tau\in\mathbb Z/a}S$, where $S=O'[[u]]$,
and $E=u^e+p$. The Frobenius matrices of the modules and their
inverses have entries in $S[1/E]$; Frobenius sends $u$ to $u^p$ and
cycles the components. After evaluation in the integral Witt ring
of the tilt with coefficients $O'$, the matrices of the two maps
are identified with the scalar extensions of the lattice maps.
The short exact sequence of free $O'$-modules splits linearly.
Consequently the two evaluated matrices have respectively a left
and a right inverse, their composite is zero, and the ranks add.
The integral evaluation is injective. One can see this on each
unramified component from the $u$-expansion: modulo the coefficient
uniformizer the map evaluates $u$ at the element
$(\pi_j\bmod p)_j$ of positive finite tilt valuation, while a
nonzero residue constant has valuation zero. The first nonzero
coefficient therefore detects the lowest term. Separation in the
coefficient uniformizer gives injectivity before reduction, with
the unramified idempotents separating the components.
In particular $GF=0$.

For each component and each of the two matrices at least one minor
of the required full rank has a coefficient in $O'^\times$.
Indeed the Cauchy--Binet identity for its evaluated one-sided inverse
expresses $1$ in the ideal of the evaluated full minors. If all
these minors in component $\tau$ were divisible by a uniformizer
$\lambda$ of $O'$, multiplication by the image of the component
idempotent would put that nonzero idempotent in $(\lambda)$.
Idempotence would then put it in every $(\lambda^n)$, contradicting
separation. Rank zero uses the empty minor $1$.

Each chosen primitive minor has only finitely many roots in the open
unit disc by Weierstrass preparation over $O'$. Thus only finitely
many points in that disc, across all components, could have deficient
rank. Away from the roots of $E$, the Frobenius intertwining identity
identifies the rank at $u=\beta$ in one component with that at
$u=\beta^p$ in its neighbor. If $0<|\alpha|<1$ were deficient,
among its $p$-th roots at most one could be a root of $E$ since
$(p,e)=1$. Choose a different root and repeat. The increasing norms
give infinitely many deficient points. At zero, the least
$u$-adic orders $d_\tau$ of the full minors are finite.
Over $L'[[u]]$ the Frobenius matrices are invertible at zero.
Change of basis preserves the ideal of the full minors, so
$d_\tau=p d_{\tau-1}$ cyclically. All $d_\tau$ are zero.

The height-one primes of the regular local ring $S$ are generated by
$\lambda$ or, up to units, irreducible distinguished polynomials,
again by preparation. At every such prime the two maps are therefore
split in their respective directions. The kernel of $F$ vanishes
by generic injectivity. For $x\in\ker G$ the rank equality and
$GF=0$ give a unique rational preimage under $F$. Local splitting
makes it integral at every height-one prime, hence integral over
$S$ since $S$ is a normal domain. This proves $\ker G=\mathrm{im}F$.
The finite module $\mathrm{coker}\,G$ is supported at the closed
point in each two-dimensional component, hence has finite length.
These are exactly the three properties of $\mathrm{IsExactTriple}$.
% END Q84T3

% BEGIN Q84T4
\paragraph{Q84T4. The exponent supplied by the monodromy estimates.}
On the common field of Q84T2 put
\[
 w=2(N-1)R,\quad h_L=(N+1)b,\quad c_1=b(R+1),\quad
 c_2=c_1+w,\quad k_*=c_2+b+h_L+w.
\]
Here $M$ is the integral module in each Breuil--Kisin datum,
$Y=(E^{-b}M)^\Delta$, $Y'=u^eY$, and $M^\Delta$ denotes the tame
invariants. Write $\mathrm{Mon}$ for its linear monodromy obstruction
in the principal-parts module. The following calculation is conditional
on the realized tame models with the indicated heights, jumps and
characters, and on three integral monodromy estimates:
\begin{enumerate}
\item Each $c\in Z_f^1(T_2)$ has a representative $y_2\in Y_2$ for
$p^w c$ with $\mathrm{Mon}(y_2)=0$.
\item Every pair $y'_2\in Y'_2$, $m_2\in M_2^\Delta$ with
$\mathrm{Mon}(y'_2+\varphi m_2)=0$ has a pair
$y'\in Y'_1$, $m\in M_1^\Delta$ with zero obstruction and
$G(y'+\varphi m)=p^{k_*}(y'_2+\varphi m_2)$.
\item For such a zero-obstruction pair, $p^{c_1}(y'+\varphi m)$
represents a class of $Z_f^1(T_1)$ under the integral period
comparison. The representation of classes is functorial for $G$,
is invariant under addition of $(\varphi-1)M^\Delta$, and two
classes represented by the same vector differ by an integral coboundary.
\end{enumerate}
Indeed $E^b-p^b=u^e\sum_{i=0}^{b-1}p^i E^{b-1-i}$. Multiplying
$y_2$ by this identity gives $p^b y_2=y'_2+m_2$, with
$m_2=E^b y_2\in M_2^\Delta$ and
$y'_2=-u^e\sum_i p^iE^{b-1-i}y_2\in Y'_2$.
Thus $y'_2+\varphi m_2=p^b y_2+(\varphi-1)m_2$. The principal-part
obstruction vanishes on the added difference, and the integral class
is multiplied by $p^b$. Apply the second and third estimates and compare
the two representatives in the target. This gives
\[
 p^dZ_f^1(T_2)\subseteq g_*Z_f^1(T_1)+B^1(T_2),\qquad
 d=c_1+k_*+b+w=6(N-1)R+(N+5+2R)b.
\]
Since $N\ge2$ and $R\ge1$, $d\le7(NR+b(N+R))$.
Multiplication by the unused power and Q84T1 give the required
constant $C_T=40$ with the original prime guard, and likewise any
$C_T\ge40$ with its stronger guard.

The formal source derives the three estimates from its
Bhatt--Scholze lattice and Fontaine period hypotheses. In the
middle estimate, its first quantitative step is the bound $p^w$ on
the cokernels of $G$ on $M,M^\Delta,Y'$. Finite length in Q84T3
does not specify this exponent. The subsequent steps bound the
obstruction by $p^{c_1}$, pull its integral principal part into
the kernel at cost $p^w$, and solve the kernel monodromy equation
at additional cost $p^{b+h_L}$. Returning to cocycles further
requires the $p^{c_1}$ bound on the comparison vector $\zeta$.
The fuller human proof of these estimates, including the initial
construction of the bounded representative, remains due. The
independent roadmap B8.42bs1d (revision \texttt{e481c74135c4})
also invokes the quantitative construction in its fourth claim.
% END Q84T4

% BEGIN Q84T5
\paragraph{Q84T5. Interpolating the two degree intervals and cancelling their common shift.}
Write $\mathcal L_i$ for the span of the Frobenius columns of $M_i$
over the integral power-series ring. Set $v=u^e$, $B=w+1$ and
$J=H+b+B$. We have $2J+2<p$: use $b<H$, $4R\le N^2$ and
$40(HN^2+N^3)<p$. Fix a tame isotypic index $c\in\mathbb Z/e$.
For $s\in A$ put $d_s=(c-s)\bmod e$ in $[0,e)$. Let $P_1,P_2$
be the monic products with roots $d_s+je$ for $0\le j<J$ and
$J\le j<2J$, respectively. Then
\[
         p^w=U_1 P_1+U_2 P_2,\qquad U_i\in\mathbb Z_p[X].
\]
To see the actual denominators, fix any root $z=d_s+je$ in the union.
In its own residue progression all nonzero differences from $z$ are
units since $2J<p$ and $p\nmid e$. For each $s'\ne s$ let $r\le R$
be the $p$-orbit length of $s'-s$. Then
$m=(d_s-d_{s'})(p^r-1)/e$ is an integer with $|m|<p^r-1$.
For another index $j'$ the difference $z-(d_{s'}+j'e)$ has
the valuation of $m+(j-j')(p^r-1)$, since both $e$ and $p^r-1$
are $p$-units. This integer is nonzero and has absolute value
$<2Jp^r<p^{r+1}$, giving valuation at most $r$. At most one value
of $j'$ in $[0,2J)$ gives positive valuation. Thus the full
Lagrange denominator at $z$ has valuation at most $(N-1)R\le w$.
Interpolate a polynomial $Q$ equal to $p^w$ on the first interval
and zero on the second. The Lagrange formula gives integral
coefficients, and division by the monic $P_i$ gives $P_2\mid Q$,
$P_1\mid p^w-Q$ over $\mathbb Z_p$. This proves the identity.

Here is the remaining module implication. Write $\mathcal H_i$ for the
high-degree submodule of the extended lattice, generated there
by scalar coefficients divisible by $u^{ep}$ in the power-series
field of coefficients. Use the following precise reductions for
the Euler connection $\nabla$: the extended lattice is
$\mathcal L_i+\mathcal H_i$ and is $\nabla$-stable. On each $c$
part, $P_1(\nabla)$ takes $\mathcal L_i$ into
$v^J\mathcal L_i+\mathcal H_i$, and $P_2(\nabla)$ takes
$v^J\mathcal L_i$ into $v^{2J}\mathcal L_i+\mathcal H_i$.
The submodule
$T=v^JG\mathcal L_1+p^B\mathcal L_2+\mathcal H_2$ is
$\nabla$-stable. These are the integral connection inputs needed here.
The tame inertia projector to $c$ preserves the integral lattices:
average with its inverse character over the cyclic group of order
$e$, a $p$-unit. The projectors sum to $1$, preserve $v^J$ and
commute with $G$. The given morphism is horizontal,
$G\nabla_1=\nabla_2G$, so its action commutes with evaluation of
every polynomial in the two connections.

Apply F84R2 with modulus $p^B$ to obtain the coarse congruence $v^J\mathcal L_2\subseteq
G\mathcal L_1+p^B\mathcal L_2$. For any $x\in\mathcal L_2$, project
$v^J x$ to $y$ in its $c$ part and apply the projected coarse
congruence. Write $y=Gz+p^Bz'$ with each term in its $c$ part.
Then $P_1(\nabla)y\in T$: on the source term use the first degree
bound and $G\mathcal H_1\subseteq\mathcal H_2$, and on the error
use stability of the extended lattice. Also $P_2(\nabla)y\in T$:
apply the second degree bound and use
$v^{2J}\mathcal L_2\subseteq v^JG\mathcal L_1+p^B\mathcal L_2$
by the same coarse congruence. The interpolation identity and
stability of $T$ give $p^w y\in T$. Sum the projectors. We obtain
\[
       v^J p^w x=v^J Gz+p^B z'+h,\qquad h\in\mathcal H_2 .
\]
The equality also puts $h$ in $\mathcal L_2$. In the Frobenius-column
basis its coefficients belong to $O'[[u]]$ and are separately
divisible by $u^{ep}$ over $L'[[u]]$. Removing their initial zero
coefficients therefore gives $h=v^ph_0$ with $h_0\in\mathcal L_2$.
Now $v^J(p^wx-Gz-v^{p-J}h_0)=p^Bz'$. In the same basis multiplication
by $v^J$ shifts the coefficient indices by $eJ$, so every coefficient
of the expression in parentheses belongs to $p^BO'$. Finally
$v^{p-J}h_0=v^J(v^{p-2J}h_0)$ belongs to
$G\mathcal L_1+p^B\mathcal L_2$ by the coarse congruence.
Consequently $p^w\mathcal L_2\subseteq G\mathcal L_1+
p^{w+1}\mathcal L_2$. Finite length of the cokernel on $M$ gives
some $p^k$ annihilating it. Applying Frobenius and taking spans
gives the same coarse annihilator on $\mathcal L$. Iterating the
last inclusion $k$ times therefore eliminates the error altogether,
proving $p^w\mathcal L_2\subseteq G\mathcal L_1$ under the stated
connection reductions and the coarse congruence.
% END Q84T5

\section{The Frobenius bound in the integral lifting argument}

% BEGIN F84R1
\paragraph{F84R1. A uniform power of the variable modulo a power of the prime.}
Let $A=O[[u]]$, let $\phi$ fix the coefficients and replace
$u$ by $u^p$, and let $I_\tau$ be ideals indexed cyclically
by $\tau\in\mathbb Z/a$ with $a\ge1$. Write
$E=u^e+p$. Suppose $1\le b$, $h+b\le p$, each
$I_\tau+(p^b)$ contains a power of $u$, and
$E^h I_{\tau+1}\subset A\phi(I_\tau)$.
Then, simultaneously at every $\tau$,
\[
               u^{e(h+b)}\in A\phi(I_\tau)+(p^b).
\]
To prove this, write each series uniquely as
$\sum_{j=0}^{p-1}u^j\phi(f_j)$ by grouping its coefficients
modulo $p$. If a series belongs to $A\phi(I)+(p^b)$ then
each $f_j$ belongs to $I+(p^b)$: for a term
$g\phi(z)$ the $j$th component is $g_j z$, and the operation
commutes with multiplication by $p^b$. Let $s_\tau$ be the
least exponent with $u^{s_\tau}\in I_\tau+(p^b)$.
In $(E-p)^{h+b-1}$ each term is divisible either by $E^h$
or by $p^b$. Multiplying by $u^{s_{\tau+1}}$ therefore puts
$u^{s_{\tau+1}+e(h+b-1)}$ in $A\phi(I_\tau)+(p^b)$.
Take its single nonzero component in the residue-class
decomposition. We get
\[
        p s_\tau\le s_{\tau+1}+e(h+b-1).
\]
At a maximal $s_\tau$ this gives
$(p-1)s_\tau\le e(h+b-1)\le e(p-1)$, hence
$s_\tau\le e$ at every index. Multiplying the preceding
monomial by the remaining nonnegative power of $u$ proves
the assertion. This argument uses no preliminary numerical
bound on the nilpotence exponent of $u$.
% END F84R1

% BEGIN F84R2
\paragraph{F84R2. The first approximation on the actual Frobenius lattice.}
Consider the integral models of Q84T3 with their original
heights $[-\eta,H]$. For $G:M_1\to M_2$ the given
map write $G_\tau$ for its component matrix in integral bases,
$A_{j,\tau}$ for the Frobenius matrices and $L_{j,\tau}=
A_{j,\tau}M_{j,\tau}$ for their column lattices in those
coordinates. The matrices satisfy
$G_\tau A_{1,\tau}=A_{2,\tau}\phi(G_{\tau-1})$,
and the height inequalities are
$E^H M_j\subset L_j$ and $E^\eta L_j\subset M_j$.
For $1\le b$ with $H+\eta+b\le p$ one has
\[
       (u^e)^{H+\eta+b} L_2\ \subset\ G L_1+p^b L_2 .
\]
Set $I_\tau=\operatorname{Ann}_A\operatorname{coker}(G_\tau)$.
The finite-length cokernel of Q84T3 implies that a power of $u$
lies in each $I_\tau$, since $u$ belongs to the maximal ideal.
The annihilator for the matrix $\phi(G_\tau)$ is exactly
$A\phi(I_\tau)$. Indeed apply the component decomposition of
F84R1 entrywise to an equation
$\phi(G_\tau)y=d\phi(x)$ for each integral $x$. The $j$th
component shows $d_j x\in\operatorname{im}G_\tau$.
Conversely, matrices of preimages of the multiples of standard
basis vectors by any $d_j\in I_\tau$ give preimages for
$\phi(d_j)$ after applying $\phi$, proving the other inclusion.

For $d\in I_\tau$ and any integral vector $z$ at the target
of the Frobenius coordinates, $E^\eta A_{2,\tau} z$ is integral.
Choose $m\in M_1$ with $G_\tau m=d E^\eta A_{2,\tau}z$.
Write $E^H m=A_{1,\tau}m'$ with $m'$ integral. Commutation with
Frobenius gives
$A_{2,\tau}\phi(G_{\tau-1})m'=
 A_{2,\tau}dE^{H+\eta}z$. Cancel $A_{2,\tau}$ over the
fraction field. The preceding annihilator identity gives
$E^{H+\eta}I_\tau\subset A\phi(I_{\tau-1})$.
Apply F84R1 with $h=H+\eta$. Multiplication by
$u^{e(H+\eta+b)}$ sends each vector $z$ into
$\phi(G_{\tau-1})M_1+p^b M_2$ in these integral coordinates.
Multiplying by $A_{2,\tau}$ and commuting $G$ back proves the
display on all components. It is a bound on a power of $u$
modulo $p^b$; the subsequent monodromy calculation is needed
to get a bound on a power of $p$.
% END F84R2

% BEGIN F84R3
\paragraph{F84R3. Passing a sharp cokernel bound to all the invariant lattices.}
In the same coordinates suppose that $p^w L_2\subset G L_1$.
Then $p^w M_2\subset G M_1$, and the same exponent works on
$E^{-\eta}M$, $M^\Delta$, $Y=(E^{-\eta}M)^\Delta$ and
$Y'=u^e Y$. Here $\Delta$ is the tame group used for the
invariants. For $z\in M_{2,\tau-1}$,
$A_{2,\tau}\phi(z)\in L_{2,\tau}$. Lift its $p^w$ multiple
to $A_{1,\tau}x$ with $x$ integral. Canceling $A_{2,\tau}$
as above gives $\phi(G_{\tau-1})x=p^w\phi(z)$.
On power series put
$\psi(\sum c_j u^j)=\sum c_{pj}u^j$. It preserves
integrality and satisfies $\psi(\phi(g)d)=g\psi(d)$ by direct
coefficient comparison. Thus $G_{\tau-1}\psi(x)=p^w z$,
giving the claim on $M$ without dividing the exponent by the
Frobenius degree. Scaling lifts by $E^{-\eta}$ gives the
next lattice. For an invariant target, average a lift over
$\Delta$; its order is a unit in $O$, and equivariance shows
that the average has the same image. Multiplication by $u^e$
then handles $Y'$. These implications isolate a single sharp
cokernel estimate on $L$ in the lifting calculation of Q84T4.
% END F84R3

\section{The global Selmer estimate for a specified cut}

% BEGIN S84G1
\paragraph{S84G1. Integral determinant and coboundaries on an arbitrary support.}
Use the data in $\mathrm{StepsNoCut}.\mathrm{GlobalStatement}$.
Write $X$ for the original points, $U$ for the tested subset,
and $\pi\in U\cap Z$. Put $O=O_i$, $K=\operatorname{Frac}O$,
$p=\ell_i$, $e=\ell_O(O/pO)$. Use the given good congruence
objects (in particular their actual integral algebra and
root cover) and put $\delta=\mathrm{del51}(\pi')$ at the
specified lift. Write $\ell$ for $O$-length in this section.
In the native normalization $\ell(O/\delta O)=e v_\delta$.
For the rank block and each character block write $C_b$
for the span of the supported equivariant cocycles of G79E7
with values in $(K/O)^{d_b}$.

First,
\[
   \ell(O/J_\pi)\le \ell(C_\rho)+\sum_{0\le j<t}\ell(C_j),
   \qquad J_\pi=\{h(\pi):h\in T_i,\ h|_{U\setminus Z}=0\}.
\]
Here is the support modification to G79E11--G79E12.
In the span of all $y_g=(\chi_1(g)^{-1}r(g)-1)E_1$ over
the matrix algebra $A$, identify two vectors if their
difference vanishes at every point above $U$. This quotient
is the image of restriction, a finite $O$-module in the
sense of finite generation. Choose finitely many $y_{g_q}$
generating it over $A$. The relations are now precisely the
rows vanishing on this restriction. As in G79E12 each solution
of the block row equations defines an equivariant map on
this image; injectivity of $K/O$ over the DVR extends it to
the ambient matrix module on the restricted points, hence
to a supported map on the original family. The Smith determinant
argument on evaluations at the $g_q$ gives the same sum of
block lengths, provided its determinant coefficients take
values in $J_\pi$ at $\pi'$.

Indeed those coefficients belong to the given root algebra
$T'$ by the single-block determinant condition and successive
formal block elimination in G79E11. Above $x\in U\setminus Z$
the absence of a $\chi_1$-line follows from $\mathrm{admG6}$:
a nonzero fixed vector in the twist is such a line and would
put $x$ in $Z$. As $E_1(x')$ has rank one, some $y_g(x')$
does not kill some vector $v$. Thus the column
$(y_{g_q}(x')v)_q$ is nonzero, and every relation annihilates
this column. The determinant polynomial vanishes on the
whole fiber. Expand a coefficient in the given free $T_i$-basis
of $T'$. For this fiber multiply the expansion by a Hecke
function supported at $x$ and nonzero there (multiply the
pairwise separating factors). The product vanishes everywhere,
so freeness and the domain property of $O$ force each original
coordinate to vanish at $x$. Thus every coordinate vanishes
on $U\setminus Z$ and has its value at $\pi$ in $J_\pi$.
Evaluation at $\pi'$ and summation prove the required assertion.
This works for every $U$ in the original quantifier.

The total loss in passing from these cocycles to cohomology
is at most $2(r+t-1)e v_\delta$. Indeed the chosen projector
has square $\delta E_1$ at $\pi'$, acts by $\delta$ on the
first line and by zero on the other blocks. Its rank there
is one, so $\delta\ne0$ by the distinguished block formula.
The calculation with its integral group-algebra expansion in
G79E7 and G81S3 applies before division: a coboundary
$g\mapsto(M_b(g)-1)v$ in $C_b$ forces $\delta^2 v=0$ for
every other block. This costs at most $2d_b\ell(O/\delta O)$
on that block. The first line has trivial ratio character
and zero coboundaries. The result holds when $\delta$ is a
nonunit, including at the two type primes.
% END S84G1

% BEGIN S84G2
\paragraph{S84G2. The original global inequality, including its local defects.}
Assume global Artin reciprocity. With the hypotheses and the
subset of S84G1, suppose that for all $w\mid p$ the original
$\mathrm{HfChar}$ and $\mathrm{LocBoundP}$ tests hold with
the stated $e_\rho,e_\chi$. Retain the source finiteness
conditions on the rank Selmer and on the indicated cyclotomic
Tate--Shafarevich lengths. Then
\[
\begin{split}
 \mathrm{vccCut}(i,U,\pi)\ \le{}&
  \#\{w\mid p\}(e_\rho+t e_\chi)+s_\pi
  +\sum_{0<j<t}\mathrm{shaLen}_{p}(j)\\
  &+v_p(h_E)+4(r+t)v_\delta
                        +t\sum_{q\in Q_i}v_\zeta(q).
\end{split}
\]
This is the native $\mathrm{GlobalStatement}$ with its
original arbitrary support and coefficient normalization.

For each $b$ and $w\mid p$, restrict $C_b$ to the decomposition
group and quotient its local span by its intersection with
the finite condition. Its image has length at most
$e e_b$, where $e_b=e_\rho$ for the rank block and $e_\chi$
otherwise. To verify this point, take the simultaneous columns
of $E_1$ over the points of the cover above $U$ and the
module spanned by their $y_g$-columns. Reassembly followed
by the supported equivariant map is well-defined, because it
kills vectors vanishing on this support. The resulting
finite list, vector, stable lattice submodule and map into
the block are a generator in the native local cut span,
by the calculation of G79E13. The coefficient isomorphism
of the local bridge identifies $O$ with the tested integers
and preserves $e$. The quotient of a span by its intersection
with the finite condition is its image modulo that condition.
Taking the kernel of the product of all these restriction
maps therefore costs at most $\#\{w\mid p\}e e_b$.

On the rank block the remaining classes belong to the
specified Selmer group. Outside $p,S_i,Q_i$ the inertia of
$r_x$ is trivial by $\mathrm{admG4}$, and that of $\chi_1$ is
trivial there by the decomposition at $\pi$. Hence $y_g$
vanishes on inertia there. The defining Selmer conditions
allow the remaining $S_i,Q_i$. The length of its image in
classes is at most $e s_\pi$.

On each character block the same unramified assertion holds.
At $S_i$, $\mathrm{admG5}$ and the tame-generator conditions
give either the integral polynomial relation, with its
evaluation at $\chi_1(\gamma_0)$ a unit, or the quadratic
inertia relation and $\zeta-1$ a unit. In the polynomial
case use that relation on $y_g$ and evaluate by the character
block: all the ratio characters are unramified here because
this place is away from $p$. Thus the unit annihilates the
cocycle value. In the quadratic case evaluation on the
tame generator puts the inertia image into $(K/O)[\zeta-1]$,
as in G79E8. Evaluation there detects the whole restriction:
it is additive on inertia, continuous with finite $p$-power
image, and kills the pro-residue-characteristic wild group.
Thus both alternatives give zero inertia restriction at $S_i$.
At $q\in Q_i$ the quadratic relation holds with
$\zeta_i(q)\ne1$. Taking the kernel of all these inertia
restrictions loses at most $e\sum_{q\in Q_i}v_\zeta(q)$
on this character block.

For $j=0$ the first-character clause of $\mathrm{HfChar}$
makes the remaining cocycles unramified also at $p$.
They factor through the ideal class group, with length at
most $e v_p(h_E)$ by G79E6. For $0<j<t$ the ratio is
$\chi_{\rm cyc}^{-j}$. The finite-character clause gives
local coboundaries above $p$. The classes unramified away
from $p$ and locally trivial above it have length at most
$e\,\mathrm{shaLen}_p(j)$ by the coefficient-coordinate
and inflation argument in G79E6. Its kernel consists of
global coboundaries, already bounded with the rank coboundaries
in S84G1. Summing all the block inequalities, applying the
determinant inequality there, and dividing by $e>0$ proves
the display: $2(r+t-1)v_\delta\le4(r+t)v_\delta$, since
$v_\delta\ge0$. The evaluation-ideal identification uses
only that the off-locus is $U\setminus Z$, as in P84D1.
% END S84G2

% BEGIN S84G3
\paragraph{S84G3. Application at the completely split type prime.}
On a good cut take an actual tested type index $p=g_k$, a
good point $\pi$, a tested good integral family $r$, and the
tested rank family $\rho$ with its semisimple decomposition
at $\pi$. Use the Galois hypotheses of C84F1--C84F6.
The point $\pi\in Z_{\rm good}$ witnesses
$\mathrm{HG}$ at this cut index; also $\pi\in Z$. The fourth
coefficient duty of C84F6 applies to exactly this $i,k,r,\rho,\pi$
and this decomposition. Its result
$\mathrm{WildIdempotentChoice}$ consumes $\pi\in Z$ again
and gives a record $S$ satisfying
$\mathrm{GoodCongruenceObjects}(D_{c}.\mathrm{nd},r,\rho,\pi,S)$
and the wild inertia condition for this same $S.g_0$.
In particular it furnishes the precise integral object in
S84G1 on the cut spectrum before the support is chosen for
its determinant calculation. The replacement of the model
family by this particular tested family is justified by the
integral conjugacy argument in C84F6.

For each $\pi$ use the support $U=S_\pi$ of P84D4, with
$\pi\in U$ and the asserted single local type tuple on all
of $U$. Assume the local hypotheses of T84L1--T84L6 for
this support, including the direct potentially crystalline
family there, the stated host, size and jump comparisons,
and the quantitative lifting input $\mathrm{admT}(C_T)$.
Also assume global Artin reciprocity and the native Selmer
and cyclotomic finiteness tests, as in S84G2. For $n=r+t$,
$a=\eta+t+1$, set $N=n+1$ and $u=\lfloor N^2/4\rfloor$.
The actual $e_{\max}=C_T(Nu+a(N+u))$ works for all
these places. T84L6 and the local bounds give
$e_\rho=2r e_{\max}$ and $e_\chi=2e_{\max}$. There are
$[E:\mathbb Q]=2f$ places above the completely split type
prime. Applying S84G2 on $S_\pi$ gives
\[
     c_\pi^*\le 4fn C_T\bigl(Nu+a(N+u)\bigr)+A_\pi
\]
with the full $A_\pi$ of P84D3. Here the arithmetic constraints
$2\le N\le n+1$, $1\le u$ and $4u\le N^2$ follow from
$n\ge3$. This holds simultaneously on the separately chosen
supports under the stated local tests. P84D5 absorbs
$v_p(N_G)$ into its $2fr^2$ reserve and sums the
multiplicities, proving the native paired estimate with
$C=C_T$. If the same direct potentiality and lifting tests
are given on the full cut spectrum, take that whole spectrum
as $U$ in S84G2 to obtain the full-spectrum $c_\pi$ bound
and use P84D3 directly.
% END S84G3

\section{Returning the estimate to the chosen family}

% BEGIN P84B1
\paragraph{P84B1. The native good-cut existence theorem.}
For every $D$ and every specified norm choice $a$, assume
$\mathrm{GaloisPreconditions}(D,a)$ and $\mathrm{Tie4}(D)$.
There exists a single cut $c$ such that
\[
  c.\mathrm{Good},\qquad
  \mathrm{NormChoiceFacts}(D_c,a_c),\qquad
  \mathrm{Tie4}(D_c),
  \qquad a_c=\mathrm{nchTr}(c,a).
\]
This is the whole assertion called $\mathrm{CutGivesAlt}$.
Indeed choose $c$ by P84C3. C84F1 and C84F2 transport the
residual and row data and give the cut preconditions and
$\mathrm{CutRec}$ there. C84F3--C84F6 now supply every
norm-choice fact on precisely that cut and transported norm.
T83C3 transports all four ties there. Thus the one choice of
$c$ witnesses the entire conjunction. P84C4 also gives its
commuting condition. No choice in this construction depends on
a subsequent representation, minor, Selmer bound or lift.
% END P84B1

% BEGIN P84B2
\paragraph{P84B2. A sufficient type-prime test for the original family.}
Work under $\mathrm{Hypotheses}_2$ with a given family
$\mathrm{Fm}$ satisfying $\mathrm{TwGood}$. At each original
member $(x,i)$ with $x>\mathrm{Fm}.x_1$, assume the Galois
preconditions for its specified norm choice and the four ties.
Fix one natural number $C_T$ for the family. Suppose that on
one good cut furnished by P84B1 at each such member the Smith
differences of P84D2, divided by $e$, obey the point bound
of P84D3 on every one of its native type-prime tests. Then
$\mathrm{LinkN}(\mathrm{Fm})$ holds. The same conclusion holds
if these differences are bounded directly by the native
$\mathrm{rhs}(C_T,\eta,i,\rho,\pi)$ in every such test.

Choose the specified cuts once for all members. At $x\le
\mathrm{Fm}.x_1$ use the cut retaining all indices and points,
so the cut family is defined everywhere. Its structural maps are
the restrictions of the original maps: its place data, scale,
parameter ledger and reference theta ranks are the original
ones, and its extra norm choice is $a_c$. In particular the
twist test restricts to this family at the original twist
threshold: every retained test for the cut uses the same
character, place and norm comparison in the uncut test.
P84B1, P84C4 and P84D3, or P84D1 in the direct-bound version,
give $\mathrm{FactsFAlt}(\mathrm{gd}\,\mathrm{Fm}\,c,C_T,C_T)$
for all $x>\mathrm{Fm}.x_1$. V84A5 applies at these common
constants and yields $\mathrm{LinkN}$ on the cut family.

For the return to the original family, fix a permitted lift
$j$ and write $d_j=L.dZ0(j)$. Use the original preconditions:
their good-row and theta-scalar clauses and localisation
uniqueness give the disjoint row count and injective theta map
of LN6. The rank clause of their row record identifies that
rank with $d_j$ at every tested twist. The integral operator of
P84C3 on the chosen projector cut, or the operator clause of
goodness for any chosen good cut, bounds the discarded defect
by LN5. Summing over localisations and weighting over places,
LN7 gives
\[
             \eta_{\rm sum}\le\eta'_{\rm sum}+d_j\log N_G .
\]
Here the normalized place weights times the valuation of the
rational integer $N_G$ sum to at most $\log N_G$ by rational
prime factorization; the same multiplicity cap $d_j$ applies
at each place before this summation. The uniform parameter bound
gives $\log N_G\le C_1 x^{C_1}L_{\rm sc}$.
For fixed positive scale exponent, enlarge the common cutoff
so $x^{C_1}\le b$ and $L_{\rm sc}\ge1$, as in LN9.
Since $f,m\ge1$ at a tested lift, the added term is bounded by
$d_j C_1 f m b L_{\rm sc}^2$. Add $C_1$ to the cut-family
constant. The other summand of the polynomial allowance is
nonnegative. The cut preserves the global cyclotomic data,
the $\mathrm{ShaFinOdd}$ guard and the series $\mathrm{vvTerm}$,
so both the summability test and its contribution are literally
the original ones. This proves the claim with constants chosen
before the member and the lift.

On $\mathrm{Des}(\mathcal C,3)$ the same argument works inside
each original $\mathrm{Closed}$ quantifier whenever the two
member hypotheses and the stated integral bound are supplied
there, with $\mathrm{TwGood}$ furnished by the design.
Use the intersection of their original ranges of scale
exponents and the maximum of their transport thresholds if
those bounds are obtained separately, as in A83G4.
This identifies the sufficient arithmetic work on that class.
P84B1 itself holds for every supplied pair $(D,a)$ at its
stated hypotheses.
% END P84B2

% BEGIN P84B3
\paragraph{P84B3. A single integral lifting estimate suffices at the type primes.}
Keep the family, twist, original preconditions, four ties and
published hypotheses of P84B2. In place of its tests on Hecke
indices, suppose one $C_T\in\mathbb N$ supplies the quantitative
potentially crystalline lattice-surjection estimate of T84L1 on
the chosen good cuts. It suffices to supply it on the degree-one
places, common fields and direct or dual lifting tests used
there. Then the original family satisfies $\mathrm{LinkN}$.

Here are the identifications in this implication. P84B1 fixes the
cut and its norm choice; the four transported ties supply the same
depth-zero witness for P84D4 and K84P1. For each good point
P84D4 gives one type locus for every point in its support.
For any native paired test, row admissibility on the cut gives
the base-change, attachment, twist and parameter comparisons of
K84P1. Thus its local compatibility holds at both lifts of each
type place. The rank-one character comparison and annihilating
inertia polynomial for K84P2--K84P4 are exactly the first and
third transported ties.

The period bridges here identify the inertia subgroups: under the
specified isomorphism of decomposition groups inertia is the kernel
of the action on the residue-field closure in both descriptions.
Complete splitting of $g_k$ gives $2f$ degree-one completions.
For a crystalline representation, restricting it further preserves
crystallinity and its labelled jumps. Indeed a crystalline period
basis remains a basis in the period tensor product; invariant
coefficients at the new field are its unramified constants.
For de Rham periods use the new field's constants and tensor the
filtration with that field, giving the jumps at the restricted
embedding. The crystalline tensor and dual comparisons pair the
dual and the original period bases perfectly and the Tate twist
translates the filtration by its one jump. These are the
Fontaine comparisons on the original identified local period
rings. Together with $\mathrm{tiePH}$ from the local bridge,
they are exactly the restriction, inertia and period inputs of
K84P2--K84P4 and T84L5. Consequently there is a direct potential
family on each of our point supports. T84L5 gives the dual
family when required by the sign and obtains both signed jump
bounds from the tested host data.

Now T84L1 uses the assumed quantitative lifting estimate.
T84L2--T84L6, with local finiteness and Bloch--Kato and Tate
duality, give the bounds for the rank block and for every
character. For S84G2 the residual full span at the tested point
is V84A3's residual matrix-span argument: reduce the semisimple
decomposition at that point, cancel its character terms and use
absolute irreducibility and density on the rank block. The
unramified, tame-generator and off-support inputs are those
verified there from the cut datum and C84F4--C84F5. The fourth
coefficient duty gives the original integral idempotent object
as in S84G3. Using global Artin reciprocity and the finiteness
guards of the paired test, S84G3 now bounds precisely the shorter
ideal of P84D5. The finite group contributes at most $2fr^2$.
So the paired condition holds with $C=C_T$ on each of these
cuts. The sum, its return to the uncut family and the uniform
constants are those of P84B2. This implication retains the
integral lifting estimate on the specified surjections; the
existence of potentially crystalline period bases by itself
supplies only the rational comparisons above.
% END P84B3

\endgroup
